\documentclass[11pt]{article}
% \renewcommand{\baselinestretch}{0.95}  % 或 0.95

% Change "review" to "final" to generate the final (sometimes called camera-ready) version.
% Change to "preprint" to generate a non-anonymous version with page numbers.
\usepackage[review]{acl}

% Standard package includes
\usepackage{times}
\usepackage{latexsym}

% For proper rendering and hyphenation of words containing Latin characters (including in bib files)
\usepackage[T1]{fontenc}
% For Vietnamese characters
% \usepackage[T5]{fontenc}
% See https://www.latex-project.org/help/documentation/encguide.pdf for other character sets

% This assumes your files are encoded as UTF8
\usepackage[utf8]{inputenc}

% This is not strictly necessary, and may be commented out,
% but it will improve the layout of the manuscript,
% and will typically save some space.
\usepackage{microtype}

% This is also not strictly necessary, and may be commented out.
% However, it will improve the aesthetics of text in
% the typewriter font.
\usepackage{inconsolata}

%Including images in your LaTeX document requires adding
%additional package(s)
\usepackage{graphicx}

% If the title and author information does not fit in the area allocated, uncomment the following
%
%\setlength\titlebox{<dim>}
%
% and set <dim> to something 5cm or larger.
% 自己添加的
\usepackage{amsmath}
\usepackage{caption}
\usepackage{booktabs}
\usepackage{multirow}
\usepackage{makecell}
\usepackage{amsmath}
\usepackage{tabularx}
\usepackage{bm}
\usepackage{enumitem}
\usepackage{booktabs} % for professional tables
\usepackage{amsmath}
%\usepackage{subfigure}
\usepackage{amssymb}
\usepackage{mathtools}
\usepackage{amsthm}
\usepackage{graphicx}   % To resize the table if needed
\usepackage{xcolor}     % For coloring cells/columns/text
\usepackage{colortbl}   % For coloring cells/columns
\usepackage{array}      % For advanced column specifications
\usepackage{pifont}     % For checkmark and cross symbols
\usepackage{adjustbox}   % 用于整体缩放
\usepackage{tabularx} % 
\usepackage{ragged2e}   % 
\usepackage{booktabs}
\usepackage{subcaption}
% --- Customization Area ---
% Define colors for the symbols and the highlight column
\definecolor{MyHighlight}{gray}{0.97} 
\definecolor{MyGreen}{RGB}{40, 167, 69}  % A nice green for checkmarks
\definecolor{MyRed}{RGB}{220, 53, 69}      % A subtle red for crosses

% Define custom commands for the symbols to make the table code cleaner
\newcommand{\cmark}{\textcolor{MyGreen}{\ding{51}}}% Checkmark
\newcommand{\xmark}{\textcolor{MyRed}{\ding{55}}}% Cross

% Define a new column type 'H' that will be highlighted
\newcolumntype{H}{>{\columncolor{MyHighlight}}c}
% --- End Customization Area ---

\title{Bloom-Eval: A Hierarchical Evaluation Benchmark \\ for Automatic Survey Generation Based on Bloom's Taxonomy}

% Author information can be set in various styles:
% For several authors from the same institution:
% \author{Author 1 \and ... \and Author n \\
%         Address line \\ ... \\ Address line}
% if the names do not fit well on one line use
%         Author 1 \\ {\bf Author 2} \\ ... \\ {\bf Author n} \\
% For authors from different institutions:
% \author{Author 1 \\ Address line \\  ... \\ Address line
%         \And  ... \And
%         Author n \\ Address line \\ ... \\ Address line}
% To start a separate ``row'' of authors use \AND, as in
% \author{Author 1 \\ Address line \\  ... \\ Address line
%         \AND
%         Author 2 \\ Address line \\ ... \\ Address line \And
%         Author 3 \\ Address line \\ ... \\ Address line}

\author{First Author \\
\author{First Author \\
  Affiliation / Address line 1 \\
  Affiliation / Address line 2 \\
  Affiliation / Address line 3 \\
  \texttt{email@domain} \\\And
  Second Author \\
  Affiliation / Address line 1 \\
  Affiliation / Address line 2 \\
  Affiliation / Address line 3 \\
  \texttt{email@domain} \\}

%\author{
%  \textbf{First Author\textsuperscript{1}},
%  \textbf{Second Author\textsuperscript{1,2}},
%  \textbf{Third T. Author\textsuperscript{1}},
%  \textbf{Fourth Author\textsuperscript{1}},
%\\
%  \textbf{Fifth Author\textsuperscript{1,2}},
%  \textbf{Sixth Author\textsuperscript{1}},
%  \textbf{Seventh Author\textsuperscript{1}},
%  \textbf{Eighth Author \textsuperscript{1,2,3,4}},
%\\
%  \textbf{Ninth Author\textsuperscript{1}},
%  \textbf{Tenth Author\textsuperscript{1}},
%  \textbf{Eleventh E. Author\textsuperscript{1,2,3,4,5}},
%  \textbf{Twelfth Author\textsuperscript{1}},
%\\
%  \textbf{Thirteenth Author\textsuperscript{3}},
%  \textbf{Fourteenth F. Author\textsuperscript{2,4}},
%  \textbf{Fifteenth Author\textsuperscript{1}},
%  \textbf{Sixteenth Author\textsuperscript{1}},
%\\
%  \textbf{Seventeenth S. Author\textsuperscript{4,5}},
%  \textbf{Eighteenth Author\textsuperscript{3,4}},
%  \textbf{Nineteenth N. Author\textsuperscript{2,5}},
%  \textbf{Twentieth Author\textsuperscript{1}}
%\\
%\\
%  \textsuperscript{1}Affiliation 1,
%  \textsuperscript{2}Affiliation 2,
%  \textsuperscript{3}Affiliation 3,
%  \textsuperscript{4}Affiliation 4,
%  \textsuperscript{5}Affiliation 5
%\\
%  \small{
%    \textbf{Correspondence:} \href{mailto:email@domain}{email@domain}
%  }
%}

\begin{document}
\maketitle
\begin{abstract}
The rapid advance of automatic survey generation (ASG) has created a critical evaluation challenge. Existing evaluation methods suffer from both cognitive dimensional simplification and methodological unreliability, primarily due to the over-reliance on the ``LLM-as-a-Judge'' approach. To bridge this gap, we establish Bloom-Eval, a six-tiered benchmark based on Bloom's Taxonomy that reliably evaluates ASG systems by prioritizing deterministic algorithms and introducing our GRADE approach for abstract abilities. Furthermore, we construct a large-scale, cross-disciplinary dataset of over 3,000 high-quality papers. Our empirical study on this benchmark reveals that while leading ASG systems are proficient format organizers, they remain unqualified knowledge integrators. This work aims to redefine ASG evaluation standards, shifting the research focus from the formal mimicry of surface structure to the cognitive deepening of intellectual content. Our method provides the ASG field with a systematic, reproducible, and theoretically grounded benchmark to guide future research.
\end{abstract}

\section{Introduction}
\label{sec:introduction}
A primary objective of artificial intelligence (AI) for science is to evolve AI from a mere information processor into a genuine knowledge generator \citep{lu2024AIScientist, li2024chain, yuan2025dolphin}. Synthesizing the vast and rapidly expanding corpus of academic literature is a critical bottleneck in achieving this goal. For this challenge, automatic survey generation (ASG) \citep{wang2024autosurvey} serves as an indispensable cognitive bridge \citep{hu2014automatic, chen2019automatic, sun2019automatic, hu2024hireview, zhu2023hierarchical}. ASG is the automated process of synthesizing relevant, disparate literature into a survey given a specific research topic. This technology transforms disjointed information into structured knowledge that fuels scientific discovery. However, as ASG has made great strides in emulating human workflows and diversifying content \citep{yan2025surveyforge, liang2025surveyx, wang2025llmmapreducev2}, its further advancement is impeded by two profound limitations in the current evaluation paradigm.

First, a cognitive science perspective reveals that existing evaluation methods suffer from severe cognitive flatness. Metrics like reference overlap rate \citep{yan2025surveyforge} and citation accuracy \citep{wang2024autosurvey} merely quantify lower-order cognitive skills like Memory and Comprehension. They fail to capture the true academic value of a survey: analytical depth, synthesis coherence, and creative originality. This lack of evaluation dimensions blurs the line between mere content organization and deep knowledge generation, misguiding the focus of ASG research. Second, an over-reliance on the ``LLM-as-a-Judge'' paradigm creates a critical transparency and reliability crisis. While demonstrating the potential of LLM-based evaluation, recent studies \citep{wang2024pandalm, zheng2023judging} also reveal that LLM judges remain prone to inherent biases when assessing subjective metrics. Given the high dependence of ASG evaluation on such subjective metrics, existing methods \citep{wang2024autosurvey, liang2025surveyx} typically use simple holistic scoring. These approaches generally employ models to directly output a final score in the absence of explicit reasoning traces, a ``black-box'' practice that fundamentally strips the evaluation of its interpretability, transparency, and auditability. This lack of transparency veils the true performance of ASG systems, hindering the ability to guide the progress of the ASG field effectively. 

To address these challenges, we introduce Bloom-Eval, a benchmark grounded in Bloom's Taxonomy \citep{anderson2001taxonomy}. This cognitive model categorizes cognitive skills into a six-level hierarchy ranging from basic recall to complex creation, providing a solid theoretical basis to distinguish mere information recall from insightful knowledge construction. Encompassing Memory, Comprehension, Application, Analysis, Evaluation, and Creation, Bloom-Eval enables fine-grained analysis of model capabilities across the full spectrum of intellectual activity. We employ a dual-constraint strategy to improve reliability and transparency. Ultimately, this work provides a diagnostic benchmark to evaluate and guide the development of ASG systems capable of high-level intelligence and creativity.

The contributions of this paper are:
\vspace{-2mm}
\begin{itemize}[itemsep=0pt, parsep=0pt]
    \item We establish Bloom-Eval, the first hierarchical evaluation benchmark for ASG, which provides a profound, six-level diagnosis of a system's cognitive abilities.
    
    \item We construct a large-scale, cross-disciplinary dataset from recent (2023-2025) top-tier publications, providing a rigorous foundation for evaluating system generalizability.
    
    \item To counter the reliability crisis of ``black-box'' LLM judges, we introduce a dual-constraint methodology that combines deterministic algorithms with our novel and transparent GRADE approach.

    \item We conduct a comprehensive empirical study that quantitatively reveals a significant and pervasive gap between the higher-order cognitive abilities of ASG systems and those of human experts. We will open-source the benchmark and evaluation scripts upon publication to foster future research.
\end{itemize}

% 其他的写作：hoang2010towards；hu2014automatic;chen2019automatic;lu2020multi；fan2018hierarchical；chang2024booookscore；cho2019towards;bosselut2018discourse;Wang2019paperrobot;balepur-etal-2023-expository
\section{Related Work}
\label{sec:related_work}
\subsection{Automatic Survey Generation}
The development of ASG is rooted in the broader landscape of scientific summarization~\cite{hoang2010towards, hu2014automatic, chen2019automatic, Wang2019paperrobot, lu2020multi, jin2020multi, xing2020automatic} and long-form text generation~\cite{fan2018hierarchical, bosselut2018discourse, cho2019towards, mao2022dyle, kryscinski2022booksum, chang2024booookscore}, where researchers have
explored diverse strategies for information aggregation. Building on these foundations, the specific task of automatic survey generation (ASG) has recently attracted significant research attention \cite{jha2015content, sun2019automatic, huang2021autosurvey, shao2024assisting}. Researchers have proposed various approaches to address distinct challenges within ASG \cite{hu2024hireview, zhu2023hierarchical}.
For example, AutoSurvey \cite{wang2024autosurvey} proposes a four-stage method for survey generation. To enhance the quality of such outputs, other systems target specific weaknesses. SurveyForge~\cite{yan2025surveyforge} learns from human patterns and enhances reference selection using a scholar navigation agent. SurveyX~\cite{liang2025surveyx} enhances content depth and diversity through deep literature preprocessing and visual generation. Addressing the separate challenge of scale, LLMxMapReduce-V2~\cite{wang2025llmmapreducev2} processes long inputs using stacked convolutional scaling layers for deep synthesis.


\subsection{Evaluation of ASG systems}
ASG evaluation increasingly adopts the ``LLM-as-a-Judge'' paradigm widely used in general generation \cite{zheng2023judging,wu2025writingbench,chen2025mlrbench}. Recent works apply this paradigm across various quality dimensions. For instance, AutoSurvey \cite{wang2024autosurvey} employs LLMs to assess aspects like citation quality and content relevance. In terms of structure, SurveyForge \cite{yan2025surveyforge} utilizes LLMs to evaluate outline rationality. SurveyX \cite{liang2025surveyx} uses semantic similarity and IoU-based reference overlap for broader evaluation. However, despite covering specific functional aspects, these metrics focus on fragmented attributes without a grounded theoretical basis. The field lacks a unified benchmark that organizes these dimensions into a cognitive hierarchy to systematically diagnose higher-order synthesis capabilities.


\section{Bloom-Eval}
\label{sec:bloom_eval}
To address the challenges of cognitive flatness and methodological unreliability, this section details the architecture of the Bloom-Eval benchmark (Figure~\ref{fig:bloom_framework}). Framework reliability is validated through human correlation studies across all tiers (Appendix~\ref{sec:appendix_validation}). We begin in Section~\ref{sec:Evaluation Methodologies} by introducing our dual-constraint strategy. This strategy utilizes two foundational methodologies (distributional analysis and the GRADE assessment approach) as the algorithmic foundation for all subsequent metrics. Following this, in Section~\ref{sec:dimensions}, we elaborate on the specific diagnostic metrics across the six cognitive tiers, demonstrating how the benchmark offers a comprehensive assessment.

\subsection{Evaluation Methodologies}
\label{sec:Evaluation Methodologies}
Our framework is built upon a dual-constraint strategy designed to maximize objectivity and minimize the risk of model hallucination. As summarized in Table~\ref{tab:metric_overview}, this strategy strictly regulates the operational boundary of the LLM based on task subjectivity. For quantifiable metrics, the model is either entirely absent (relying on deterministic algorithms) or confined to specific functional tools (e.g., extraction, matching); for abstract cognitive dimensions, it operates solely as a constrained judge via the transparent GRADE approach. This principled design is the key to our framework's reliability.

To implement this strategy, we integrate standard set-comparison metrics (i.e., F1-score) with two specialized foundational methodologies tailored for distinct assessment types (implementation details in Appendix~\ref{sec:detail_Evaluation_Methodologies}):

\begin{itemize}[leftmargin=0.2cm, itemsep=0pt, parsep=0pt]
    \item \textbf{Distributional Similarity (DS):} Standard overlap metrics (e.g., Jaccard) merely check for content presence, failing to detect focus misalignment. We propose DS, a composite metric fusing Jensen-Shannon Divergence, Hellinger Distance, and Total Variation Distance, to diagnose whether the model's focus distribution aligns with expert consensus.

    \item \textbf{Generative Rubric Adaptive Differential Evaluation (GRADE):} To assess abstract abilities, we introduce GRADE approach, inspired by~\cite{du2025deepresearch}. Unlike conventional ``LLM-as-a-Judge'' approaches that produce a single, opaque score, GRADE enforces a transparent process through a two-stage method. First, it externalizes criteria by guiding the Judge LLM to deconstruct the task into explicit, weighted rubrics. Second, it performs justified differential scoring, where the LLM compares the generated survey against a human expert reference. This step requires the model to output not just a numerical score but also a detailed textual analysis justifying the score difference based on the criteria. This ensures judgments are auditable, interpretable, and anchored to verifiable standards.
\end{itemize}


\begin{table}[h!]
\centering
\caption{Overview of Bloom-Eval's metrics, methodologies, and the corresponding role of the LLM.
\textbf{Method} abbreviations: 
\textit{E.} (LLM-based Extraction), 
\textit{M.} (LLM-based Matching), 
\textit{Alg.} (Algorithm). 
\textbf{LLM Role} definitions: 
\textit{Tool} for specific tasks (E./M.), 
\textit{Constrained Judge} for scoring via GRADE approach.}
\label{tab:metric_overview}

\footnotesize             
\setlength{\tabcolsep}{2pt}  % 列间距
\renewcommand{\arraystretch}{1} % 行

\begin{tabularx}{\columnwidth}{l l l c} 
\toprule
\textbf{Level} & \textbf{Metric} & \textbf{Method} & \textbf{LLM Role} \\ 
\midrule
\multirow{3}{*}{Memory} & EFid & E. + M. + Alg. &   Tool \\
 & HIRC & Alg. & None \\
 & FCons & E. + M. + Alg. & Tool \\
\midrule
\multirow{4}{*}{Comprehension} & OTC & E. + M. + Alg. & Tool \\
 & CF & E. + M. + Alg. & Tool \\
 & TFSim & Alg. + M. & Tool \\
 & TBal & Alg. & None \\
\midrule
\multirow{3}{*}{Application} & FMI & Alg. & None \\
 & DSI & Alg. & None \\
 & FAP & GRADE & Constrained Judge \\
\midrule
\multirow{3}{*}{Analysis} & STS & Alg. & None \\
 & SCons & Alg. & None \\
 & SCS & M. + Alg. & Tool \\
\midrule
Evaluation & CAA & E. + M. + Alg. & Tool \\
\midrule
\multirow{2}{*}{Creation} & FNov & GRADE & Constrained Judge \\
 & ROQ & GRADE & Constrained Judge \\
\bottomrule
\end{tabularx}
\end{table}

\subsection{Hierarchical Evaluation Dimensions}
\label{sec:dimensions}

\begin{figure*}[!t]
    \centering
    \setlength{\abovecaptionskip}{0.2pt} 
    \includegraphics[width=\textwidth]{latex/bloom-eval_10_05.png}
    \caption{An overview of the Bloom-Eval framework, detailing (a) the data collection pipeline and (b) the six-tiered hierarchical evaluation across the cognitive levels of Memory, Comprehension, Application, Analysis, Evaluation, and Creation.}
    \label{fig:bloom_framework}
\end{figure*}

This section details the diagnostic metrics for each cognitive tier, with mathematical formulations in Appendix~\ref{sec:appendix_hm} and theoretical mappings to Bloom's Taxonomy in Appendix~\ref{sec:appendix_mapping}.

\paragraph{Level 1: Memory.}
\label{sec:l1_memory}
In Bloom's Taxonomy, Memory refers to the ability to recall information, encompassing facts, concepts, definitions, and terminology. Mapped to ASG, this tier mandates that the system accurately reproduce the field's foundational elements: core concepts and terms, pivotal literature, and factual statements. We diagnose these retrieval capabilities via three metrics:
% 修改成 To..., EFid...的句式 突出insight, 添加与memory的联系
\begin{itemize}[leftmargin=0.2cm, itemsep=0pt, parsep=0pt]
    \item \textbf{Entity Fidelity (EFid):} To diagnose the integrity of domain entity memory, EFid combines the F1-score and DS score of shared entities to measure both the scope of entity coverage and the consistency of relative emphasis.

    \item \textbf{High-Impact Reference Coverage (HIRC):} To diagnose foundational literature memory, HIRC uses the F1-score to verify the inclusion of seminal works (citations $> 50$) representing the domain's intellectual heritage.

    \item \textbf{Factual Consistency (FCons):} To diagnose factual precision, FCons uses the F1-score to measure whether generated claims are logically entailed by expert ground truth via semantic matching, strictly distinguishing faithful recall from hallucination.
\end{itemize}

\paragraph{Level 2: Comprehension.}
\label{sec:l2_comprehension}
In Bloom's Taxonomy, Comprehension refers to the ability to explain and summarize. Mapped to ASG, this requires the model to summarize content into necessary sub-headings, explain cited references, and synthesize the field's research landscape and distributional hotspots. We design four metrics to capture these capabilities, addressing potential failures in topic classification, interpretative faithfulness, and thematic synthesis:

\begin{itemize}[leftmargin=0.2cm, itemsep=0pt, parsep=0pt]
    \item \textbf{Outline Topic Coverage (OTC):} To diagnose the ability to classify and identify necessary sub-topics, OTC evaluates the generated survey's coverage of the field's main categories via semantic matching of outline headings.

    \item \textbf{Citation Faithfulness (CF):} To assess the interpretative accuracy of in-text citations, CF measures whether a statement correctly paraphrases and summarizes the cited paper. We report F1-score following~\cite{wang2024autosurvey}.

    \item \textbf{Thematic Focus Similarity (TFSim):} To assess the alignment of thematic understanding, TFSim uses BERTopic~\citep{grootendorst2022bertopic} to categorize references and calculates the F1-score and DS score to quantify the consistency of the generated research focus relative to experts.

    \item \textbf{Thematic Balance (TBal):} To measure the distributional evenness of the model's research attention, TBal quantifies the equality of cited references across topics. It uses the Gini Coefficient \cite{gini1912variabilita} to detect disproportionate fixation, ensuring the model allocates attention evenly rather than exhibiting bias toward specific sub-fields.
\end{itemize}

\paragraph{Level 3: Application.}
\label{sec:l3_application}
In Bloom's Taxonomy, Application refers to applying learned knowledge to practice. Mapped to ASG, this entails applying standard academic survey frameworks, specifically necessitating basic formatting correctness, structural completeness, and framework application capabilities. We capture these capabilities via:

\begin{itemize}[leftmargin=0.2cm, itemsep=0pt, parsep=0pt]
\item \textbf{Formatting Integrity (FMI):} To diagnose adherence to citation conventions, FMI assesses the set consistency between in-text markers and the bibliography to verify the strict execution of academic referencing standards.

\item \textbf{Document Structure Integrity (DSI):} To diagnose the application of standard academic templates, DSI evaluates the presence of essential structural components (e.g., Abstract, Introduction) to verify compliance with canonical publication norms.

\item \textbf{Framework Application (FAP):} To diagnose the application of organizing paradigms, FAP uses the GRADE approach to evaluate the consistent implementation of a recognized taxonomy (e.g., chronological), assessing the transfer of abstract organizational logic to the specific survey topic.
\end{itemize}


\paragraph{Level 4: Analysis.}
\label{sec:l4_analysis}
In Bloom's Taxonomy, Analysis is the ability to break down information and identify structures or relationships. Mapped to ASG, Analysis implies ensuring an accurate writing structure, characterized by correct hierarchical logic, appropriate granularity, and structural clarity. Treating the survey's outline as a proxy for its knowledge structure, we propose the following metrics to capture these capabilities:

\begin{itemize}[leftmargin=0.2cm, itemsep=0pt, parsep=0pt]

\item \textbf{Semantic Tree Similarity (STS):} To diagnose the soundness of hierarchical logic, STS quantifies the structural and semantic similarity between generated and expert outlines (modeled as semantic trees) via edit distance, verifying consistency in relational logic.

\item \textbf{Shape Consistency (SCons):} To diagnose the appropriateness of analytical granularity, SCons compares the structural complexity (depth and breadth) of the generated and expert outlines, ensuring the analysis avoids becoming superficially flat or overly narrow.

\item \textbf{Structural Clarity Score (SCS):} To diagnose the precision of deconstruction, SCS strictly penalizes semantically redundant sections appearing in disparate branches, ensuring the clear separation of distinct concepts without confusing repetition.
\end{itemize}

\paragraph{Level 5: Evaluation.}
\label{sec:l5_evaluation}
In Bloom's Taxonomy, Evaluation is defined as the ability to critically judge the value of information. Mapped to ASG, this entails critiquing the conclusions and limitations of existing literature. We design the following metric to capture this capability:
\begin{itemize}[leftmargin=0.2cm, itemsep=0pt, parsep=0pt] 
\item \textbf{Critical Analysis Alignment (CAA):}
To diagnose the quality of evaluative reasoning, CAA measures the F1-score overlap of explicit value judgments (e.g., method flaws) between generated and expert surveys, verifying if critical insights are both accurate and exhaustive.
\end{itemize}

\paragraph{Level 6: Creation.}
\label{sec:l6_creation}
Representing the highest order of cognition, Creation emphasizes innovation, requiring the capacity to propose new viewpoints, models, or works. Mapped to ASG, this tier specifically mandates framework novelty (reorganizing the field with fresh perspectives) and future work discovery (identifying potential research directions). We evaluate two key capabilities:

\begin{itemize}[leftmargin=0.2cm, itemsep=0pt, parsep=0pt] 
\item \textbf{Framework Novelty (FNov):} To diagnose the creation of new conceptual models, FNov uses GRADE to evaluate the originality and insightfulness of the proposed framework, assessing whether it reorganizes the research field with a fresh interpretive lens rather than merely replicating existing schemas.
\item \textbf{Research Outlook Quality (ROQ):} To diagnose the discovery of future work, ROQ measures the foresight and strategic value of the research directions proposed by the system, typically found in the conclusion or future work sections.
\end{itemize}


%修改前
% \section{Bloom-Eval}


% % (e.g., Tree Edit Distance for structural analysis) or confines the role of LLMs to highly specific, verifiable sub-tasks (e.g., entity extraction, semantic matching). For abstract high-order cognitive metrics, we introduce the "Generative Rubric Adaptive Differential Evaluation (GRADE)" protocol. 

% % This protocol first guides a Judge LLM to dynamically generate specific, weighted scoring criteria for a given task and then conducts comparative evaluations based on these explicit rubrics, thereby significantly enhancing the transparency and reliability of high-order cognitive assessment.


% To address the challenges of cognitive flatness and methodological unreliability outlined previously, this section details the architecture and principles of the Bloom-Eval framework. We begin in Section~\ref{sec:Evaluation Methodologies} by defining our three foundational methodologies (set-based comparison, distributional analysis, and the GRADE assessment approach), which serve as the building blocks for all subsequent metrics. Following this, in Section~\ref{sec:dimensions}, we elaborate on the specific metrics designed for each of the six cognitive tiers, from Memory to Creation, to demonstrate how the framework offers a comprehensive and diagnostic assessment of a model's capabilities.
% % To enhance reliability, Bloom-Eval employs a dual-constraint strategy. For quantifiable, low-order cognitive metrics, it restricts the use of LLMs to well-defined sub-tasks like information extraction and classification, or replaces their judgment entirely with deterministic algorithms such as Tree Edit Distance. For abstract, high-order metrics, it introduces the GRADE protocol, which anchors LLM judgments by dynamically generating scoring criteria and performing a differential evaluation, thereby improving assessment reliability.


% \subsection{Hierarchical Evaluation Dimensions}
% \label{sec:dimensions}

% \begin{figure*}[!t]
%     \centering
%     \setlength{\abovecaptionskip}{0.2pt} 
%     \includegraphics[width=\textwidth]{latex/bloom-eval_10_05.png}
%     \caption{An overview of the Bloom-Eval framework, detailing (a) the data collection pipeline and (b) the six-tiered hierarchical evaluation across the cognitive levels of Memory, Comprehension, Application, Analysis, Evaluation, and Creation.}
%     \label{fig:bloom_framework}
% \end{figure*}

% % \paragraph{Level 2: Comprehension.}
% % \label{sec:l2_comprehension}
% % This tier evaluates the model's ability to understand information by interpreting, paraphrasing, and summarizing it.

% % % The ability of a model to summarize the high-level structure of a topic is quantified by Outline Topic Coverage, while the capacity of the model to accurately paraphrase concepts from source material is measured by Citation Content Consistency. Beyond these specific skills, a model with strong understanding should produce a survey that is thematically congruent with one written by an expert. This holistic alignment is evaluated through a bibliographic analysis; Thematic Focus Similarity and Thematic Balance assess the thematic distribution of the cited references to measure congruence in research coverage and focus.

% % \begin{itemize}[leftmargin=0.2cm, itemsep=0pt, parsep=0pt]
% %     \item \textbf{Outline Topic Coverage (OTC):}     
% %     OTC evaluates whether the generated survey's outline identifies and covers the main topics within the field. We extract all section headings from the generated (\(O_{\text{gen}}\)) and expert (\(O_{\text{gold}}\)) outlines into two topic sets, use an LLM for semantic matching between them, and then compute the F1 as in \eqref{eq:1}.

% %     \item \textbf{Citation Faithfulness (CF):} 
% %     CF measures the faithfulness of in-text citations by assessing whether a statement accurately summarizes its cited source. We adopt the Citation Recall and Citation Precision metrics from~\cite{wang2024autosurvey} and report the corresponding F1.

% %     \item \textbf{Thematic Focus Similarity (TFSim):} 
% %     TFSim assesses the alignment of research focus between the generated and expert surveys by analyzing their bibliographies. We use BERTopic \citep{grootendorst2022bertopic} for theme clustering and report both the F1 on theme overlap and the DS score on their distribution.

    

% %     \item \textbf{Thematic Balance (TBal):} TBal quantifies the balance of topic coverage within a single survey. We apply the standard Gini Coefficient (G) \cite{gini1912variabilita} to capture the distributional inequality. The final metric is defined as \(1 - G\):
% %     {
% %     \setlength{\abovedisplayskip}{6pt}%  % 缩小公式上方间距
% %     \setlength{\belowdisplayskip}{6pt}%  % 缩小公式下方间距
% %     \begin{equation}
% %     G = \frac{\sum_{i=1}^{n} \sum_{j=1}^{n} |x_i - x_j|}{2n \sum_{i=1}^{n} x_i} ,
% %     \end{equation}
% %     }
% %     where $x_i$ is the number of references in theme $i$. A higher TBal score thus reflects a more comprehensive research focus.

    
% % \end{itemize}

% % \paragraph{Level 3: Application.}
% % \label{sec:l3_application}
% % This tier assesses the ability to apply learned knowledge and adhere to established academic conventions in the writing task.

% % \begin{itemize}[leftmargin=0.2cm, itemsep=0pt, parsep=0pt]
% %     \item \textbf{Formatting Integrity (FMI):} 
% %     This metric assesses the consistency between in-text citations and the bibliography to detect errors like ``dangling citations''. It is calculated as the Jaccard similarity between the set of citation identifiers found in the text \(C_{\text{cited}}\) and those defined in the bibliography \(R_{\text{defined}}\):
% %     {
% %     \setlength{\abovedisplayskip}{6pt}%  % 缩小公式上方间距
% %     \setlength{\belowdisplayskip}{6pt}%  % 缩小公式下方间距
% %     \begin{equation}
% %     \text{FMI} = \frac{|C_{\text{cited}} \cap R_{\text{defined}}|}{|C_{\text{cited}} \cup R_{\text{defined}}|}.
% %     \end{equation}
% %     }
% %     \item \textbf{Document Structure Integrity (DSI):} This checklist-based metric measures the presence of pre-defined, essential sections $S_{\text{req}}$ (e.g., Abstract, Introduction, Conclusion, References) within the generated document $S_{\text{found}}$. The score is calculated as:
% %     {
% %     \setlength{\abovedisplayskip}{6pt}%  % 缩小公式上方间距
% %     \setlength{\belowdisplayskip}{6pt}%  % 缩小公式下方间距
% %     \begin{equation}
% %         \text{DSI} = \frac{\sum_{s \in S_{\text{req}}} \mathbb{I}(s \in S_{\text{found}})}{|S_{\text{req}}|},
% %     \end{equation}
% %     }
% %     where \(\mathbb{I}(\cdot)\) is an indicator function.

% %     \item \textbf{Framework Application (FAP):} This metric evaluates the model's ability to select and consistently apply a recognized and appropriate organizing paradigm (e.g., chronological, methodological). The final score is calculated according to the GRADE approach detailed in Section~\ref{sec:grade}.
% % \end{itemize}


% % \paragraph{Level 4: Analysis.}
% % \label{sec:l4_analysis}
% % This tier evaluates a model's analysis skills by quantitatively assessing its ability to deconstruct and organize information into a coherent structure. We treat the survey's outline as a proxy for its knowledge structure.

% % \begin{itemize}[leftmargin=0.2cm, itemsep=0pt, parsep=0pt]
% % \item \textbf{Semantic Tree Similarity (STS):}
% % This metric quantifies the hierarchical and semantic similarity between the generated outline (Tree A) and the expert outline (Tree B). We parse both outlines into tree structures and compute their similarity using the Tree Edit Distance (TED) algorithm. The cost for node insertion and deletion is 1, while the update cost for a node (i.e., a section heading) is its semantic distance, defined as \(1 - \cos(\theta)\) based on vector embedding similarity. The final normalized score is calculated as:
% % {
% % \setlength{\abovedisplayskip}{6pt}%  % 缩小公式上方间距
% % \setlength{\belowdisplayskip}{6pt}%  % 缩小公式下方间距
% % \begin{equation}
% % \text{Similarity} = 1 - \frac{\text{TED}(A, B)}{|A| + |B|},
% % \end{equation}
% % }
% % where \(|A|\) and \(|B|\) are the number of nodes in each tree. A score approaching 1 indicates high structural and semantic congruence.

% % \item \textbf{Shape Consistency (SCons):}
% % This diagnostic metric compares the structural style of two outlines by quantifying the similarity of their depth and breadth. It is calculated as the geometric mean of two ratios: \textbf{Depth Consistency (DC)} and \textbf{Breadth Consistency (BC)}. Let the depths of the two outlines be $d_1$ and $d_2$, and their total node counts be $n_1$ and $n_2$. The consistency ratios are defined as the smaller value divided by the larger value for each dimension:
% % {
% % \setlength{\abovedisplayskip}{6pt}%
% % \setlength{\belowdisplayskip}{6pt}%
% % \begin{equation}
% %     \text{DC} = \frac{\min(d_1, d_2)}{\max(d_1, d_2)} \quad \text{and} \quad \text{BC} = \frac{\min(n_1, n_2)}{\max(n_1, n_2)}.
% % \end{equation}
% % }
% % The final Shape Consistency score is then computed as:
% % {
% % \setlength{\abovedisplayskip}{6pt}%
% % \setlength{\belowdisplayskip}{6pt}%
% % \begin{equation}
% %     \text{SCons} = \sqrt{\text{DC} \times \text{BC}}.
% % \end{equation}
% % }

% % \item \textbf{Structural Clarity Score (SCS):}
% % This metric assesses an outline's logical coherence by quantifying improper content duplication. We use an LLM to identify semantically equivalent heading pairs that do not share an immediate parent, a rule designed to prevent misidentifying normal sibling sections as redundant. These pairs are counted as redundant \(N_{\text{redundant}}\). The clarity score is then calculated as:
% % {
% % \setlength{\abovedisplayskip}{6pt}%  % 缩小公式上方间距
% % \setlength{\belowdisplayskip}{6pt}%  % 缩小公式下方间距
% % \begin{equation}
% %     \text{SCS} = 1 - \frac{N_{\text{redundant}}}{N_{\text{total}}},
% % \end{equation}
% % }
% % where \(N_{\text{total}}\) is the total number of sections. A score near 1 suggests a clear structure with minimal content overlap.
% % \end{itemize}



% % \paragraph{Level 5: Evaluation.}
% % \label{sec:l5_evaluation}
% % This tier evaluates the model's critical thinking by quantifying its ability to judge the value and reliability of existing work in alignment with an expert.

% % \begin{itemize}[leftmargin=0.2cm, itemsep=0pt, parsep=0pt]
% % \item \textbf{Critical Analysis Alignment (CAA):} CAA quantifies the alignment between the AI's and an expert's critical statements (e.g., evaluations, limitations) by extracting and semantically matching them, then calculating the F1. In this context, Precision reflects the accuracy of the AI's insights, while Recall measures the comprehensiveness of the system’s analysis.
% % \end{itemize}


% % \paragraph{Level 6: Creation.}
% % \label{sec:l6_creation}
% % This tier assesses the system's ability to synthesize existing knowledge into original insights. Using the GRADE approach, we evaluate two key creative abilities:
% % \begin{itemize}[leftmargin=0.2cm, itemsep=0pt, parsep=0pt]
% %     \item \textbf{Framework Novelty (FNov):} This measures the originality and insightfulness of any new conceptual framework proposed by the system to reorganize or reinterpret the research field. The GRADE approach is prompted to specifically focus on the originality and explanatory power of the proposed framework.
% %     \item \textbf{Research Outlook Quality (ROQ):} This measures the quality and foresight of the new research questions and future directions proposed by the system, typically found in the conclusion or future work sections.
% % \end{itemize}







\section{Experiments}

\subsection{Benchmark Construction and Topic Selection}
\label{sec:bench_topic}
To conduct a comprehensive, equitable, and rigorous evaluation of ASG systems, we first established a benchmark characterized by its alignment with the scientific frontier, high quality, and disciplinary diversity.
\begin{figure}[h]
\centering
\includegraphics[width=0.5\textwidth]{latex/Distribution_of_annual_reviews.pdf}
\caption{Distribution of Annual Reviews articles across 51 journal series (2023-2025).}
\label{fig:distribution_of_annual_reviews}
\end{figure}

\paragraph{Benchmark Corpus Construction.}
We constructed an authoritative corpus of 3,506 manually verified survey papers as the expert ground truth, curated based on three core principles: quality, scope, and timeliness. To ensure quality, all papers were selected from top-tier academic conferences (e.g., ICLR, ACL) and prestigious journal series (e.g., Annual Reviews), strictly excluding content from preprint servers like arXiv that has not undergone rigorous peer review. To achieve a broad scope for cross-domain evaluation, the corpus spans both Computer Science (CS) and a wide array of General Science disciplines. The disciplinary breakdown of the General Science portion, sourced from the Annual Reviews series, is detailed in Figure~\ref{fig:distribution_of_annual_reviews}. Finally, all papers were published between 2023-2025 to ensure timeliness and minimize data leakage risks. 
% This ensures the relevance of our evaluation topics to the current research landscape and minimizes potential data leakage from the pre-training corpora of large language models.

A comparative analysis presented in Figure~\ref{fig:data_analysis_1} demonstrates the superiority of our benchmark over prior works \citep{wang2024autosurvey, yan2025surveyforge, liang2025surveyx, wang2025llmmapreducev2}. 

The analysis reveals two critical advantages of Bloom-Eval. First, regarding scope, the figure shows that while existing benchmarks are almost exclusively confined to a few computer science sub-fields, our benchmark provides balanced coverage across all 14 domains, including comprehensive collections in the general sciences. Second, in terms of scale, our dataset of 3,506 papers is substantially larger than its predecessors.

Table~\ref{tab:coverage_matrix} highlights the contrast in data source quality and scale. While prior works \citep{guo2025sgsimeval} rely almost exclusively on arXiv preprints, our benchmark is built upon a diverse set of 60 top-tier, peer-reviewed venues, including 9 major CS conferences and 51 journals from the Annual Reviews series. This distinction in data sources underscores the superior authoritativeness and reliability of our benchmark.

% To guarantee the authoritativeness and reliability of the benchmark, papers accepted only as abstracts or those published in informal venues like workshops are excluded.

% This extensive corpus provides a solid foundation for our evaluation. To visually underscore its advantages in both scale and scope, Figure~\ref{fig:data_analysis_1} provides a comparative analysis against several existing benchmarks \citep{wang2024autosurvey, yan2025surveyforge, liang2025surveyx, wang2025llmmapreducev2}. As the figure illustrates, our dataset is not only nearly 14 times larger than previous benchmarks but also spans 14 distinct topic categories from various subfields of computer science and other disciplines. The figure highlights a critical distinction: while prior benchmarks are overwhelmingly concentrated in a few sub-fields of computer science, our proposed benchmark offers substantially broader and more balanced coverage across all domains.

\begin{figure}[h]
    \centering
    \includegraphics[width=0.5\textwidth]{latex/domain.png}
    \caption{Comparative analysis of benchmark coverage across domain categories.}
    \label{fig:data_analysis_1}
\end{figure}

\begin{table*}[t]
\centering
\caption{Data source composition and scale of survey generation benchmarks.}
\label{tab:coverage_matrix}

\footnotesize             
\setlength{\tabcolsep}{2pt}  % 列间距
\renewcommand{\arraystretch}{1} % 行

\begin{tabular}{l ccccc H}
\toprule
\textbf{Venue} & \textbf{AutoSurvey} & \begin{tabular}{@{}c@{}}\textbf{SurveyForge} \\ \textbf{(SurveyBench)}\end{tabular} & \begin{tabular}{@{}c@{}}\textbf{LLM×MapReduce-V2} \\ \textbf{(SurveyEval)}\end{tabular} & \textbf{SGSimEval} & \textbf{SurveyX} & \textbf{Bloom-Eval (Ours)} \\
\midrule
arXiv                          & \cmark & \cmark & \cmark & \cmark & \cmark & \xmark \\
AAAI                           & \xmark & \xmark & \xmark & \xmark & \xmark & \cmark \\
ACL                            & \xmark & \xmark & \xmark & \xmark & \xmark & \cmark \\
EMNLP                          & \xmark & \xmark & \xmark & \xmark & \xmark & \cmark \\
ICLR                           & \xmark & \xmark & \xmark & \xmark & \xmark & \cmark \\
IJCAI                          & \xmark & \xmark & \xmark & \xmark & \xmark & \cmark \\
IJCV                           & \xmark & \xmark & \xmark & \xmark & \xmark & \cmark \\
JMLR                           & \xmark & \xmark & \xmark & \xmark & \xmark & \cmark \\
NIPS                           & \xmark & \xmark & \xmark & \xmark & \xmark & \cmark \\
TPAMI                          & \xmark & \xmark & \xmark & \xmark & \xmark & \cmark \\
\textit{Annual Reviews (51)}   & \xmark & \xmark & \xmark & \xmark & \xmark & \cmark \\
\midrule
\textbf{Total Venues}          & 1      & 1      & 1      & 1      & 1      & \textbf{60} \\
\textbf{Total Papers}          & 20     & 100    & 384    & 80     & 20     & \textbf{3,506} \\
\bottomrule
\end{tabular}
\end{table*}
\paragraph{Topic Selection.} As a full-scale evaluation was impractical, we selected 20 representative topics based on impact. We chose the 10 most-cited papers from both our ``Computer Science'' and ``General Science'' categories. The titles of these 20 papers became our evaluation topics, guaranteeing a test set that is authoritative, diverse, and timely. These 20 topics represent specific, well-defined academic questions rather than broad concepts, providing precise targets for assessing the performance of ASG systems.


% In summary, through stringent source selection, large-scale manual collection, and a strong emphasis on diversity, our dataset presents a more challenging benchmark than existing alternatives, providing a solid foundation for the evaluation and advancement of next-generation ASG systems.




\subsection{Experimental Setup}

\paragraph{Implementation Details.} For reproducibility, we use the gpt-5-mini model with a temperature of $0.0$ for internal tasks like content extraction, matching and scoring (see Appendix~\ref{sec:appendix_prompts} for the exact prompts used).
To ensure comparability with prior work, we follow the methodology of Autosurvey~\cite{wang2024autosurvey} for specific modules, using gpt-4o for citation metrics and nomic-ai/nomic-embed-text-v1 for text vectorization.

\paragraph{Evaluated ASG Methods.}
For a fair comparison, all ASG methods were evaluated using their official default settings from public codebases \citep{wang2024autosurvey, yan2025surveyforge} or publicly available online services \citep{liang2025surveyx, wang2025llmmapreducev2}. Our benchmark includes topics from both Computer Science (CS) and general science domains to test system robustness and generalization. We acknowledge that some systems, such as AutoSurvey~\cite{wang2024autosurvey}, are designed specifically for the CS domain. Their inclusion here is intended not for direct comparison on an unfair basis, but to systematically measure their out-of-domain performance and understand the capability boundaries of current tools. Therefore, the performance of such systems on non-CS topics should be interpreted in this context. This approach provides a valuable baseline and reference for exploring the cross-domain adaptability of ASG systems.


\subsection{Main Results}

\begin{figure}[t]
    \centering
    \includegraphics[width=0.4\textwidth]{latex/all_levels.png}
    \caption{High-level performance of ASG systems across the six cognitive tiers of Bloom-Eval. Each axis score is the arithmetic mean of its tier's constituent metrics. For composite metrics (P/R/F1), the F1 score is used; for metrics with both F1 and DS, their average is taken.}
    \label{fig:all_levels}
\end{figure}

% \begin{figure}[h]
%     \centering
%     \includegraphics[width=0.4\textwidth]{latex/all_levels.png}
%     \caption{Performance of ASG Systems on the Bloom-Eval Framework.}
%     \label{fig:all_levels}
% \end{figure}

\begin{table*}[t]
    \centering
    \caption{
        Comprehensive performance of ASG systems across the six cognitive tiers of Bloom-Eval. The Human column provides reference scores where a dash (-) indicates a definitional value: 1.0 for gold-standard comparisons (e.g., F1-score) and 0.5 as the reference point for human parity in the GRADE metrics. Other numerical scores are empirical measurements of the expert-written surveys. Detailed results in Appendix~\ref{sec:appendix_detailed_results}.
    }
    \label{tab:main_tab}
    \footnotesize
    \setlength{\tabcolsep}{2pt} 
    \renewcommand{\arraystretch}{1} % 稍微增加行高以提高可读性
    \begin{tabular}{l l | c c | c c | c}
        \toprule
        % --- Table Header ---
        \multicolumn{2}{c|}{} & \multicolumn{2}{c|}{\textbf{Domain-Specific Systems}} & \multicolumn{2}{c|}{\textbf{General Domain Systems}} & \multirow{2}{*}{\textbf{Human}} \\
        \cmidrule(lr){3-4} \cmidrule(lr){5-6}
        
            \textbf{Level} & \textbf{Metric} & \textbf{AutoSurvey} & \textbf{SurveyForge} & \makecell{\textbf{LLMxMap} \\ \textbf{Reduce-V2}} & \textbf{SurveyX} & \\
        \midrule
        
        % --- Level 1: Memory 已核对---
        \multirow{4}{*}{\textbf{Memory}} 
        & Entity Fidelity (F1/DS)      & 0.10 / 0.81 & 0.13 / 0.73 & 0.12 / 0.69 & 0.11 / 0.78 & - \\
        & \makecell[l]{High-Impact Reference Coverage (F1)}         & 0.03 & 0.07 & 0.02 & 0.04 & - \\
        & Factual Consistency (F1)     & 0.07 & 0.06 & 0.09 & 0.07 & - \\
        \cmidrule(lr){2-7}
        & \textit{Avg. Score} & 0.19 & 0.19 & 0.17 & 0.19 & - \\
        \midrule

        % --- Level 2: Comprehension ---
        \multirow{5}{*}{\textbf{Comprehension}} 
        & Outline Topic Coverage (F1)          & 0.40 & 0.61 & 0.53 & 0.56 & - \\
        & Citation Faithfulness (F1)           & 0.03 & 0.27 & 0.29 & 0.37 & 0.47 \\
        & \makecell[l]{Thematic Focus Similarity (F1/DS)}     & 0.43 / 0.73 & 0.53 / 0.78 & 0.41 / 0.67 & 0.52 / 0.81 & - \\
        & Thematic Balance              & 0.72 & 0.70 & 0.68 & 0.68 & 0.70 \\
        \cmidrule(lr){2-7}
        & \textit{Avg. Score} & 0.43 & 0.56 & 0.51 & 0.57 & - \\
        \midrule

        % --- Level 3: Application ---
        \multirow{4}{*}{\textbf{Application}} 
        & Formatting Integrity               & 1.00 & 1.00 & 0.95 & 1.00 & 0.99 \\
        & \makecell[l]{Document Structure Integrity}               & 0.74 & 0.75 & 0.69 & 1.00 & 1.00 \\
        & Framework Application              & 0.38 & 0.42 & 0.52 & 0.41 & - \\
        \cmidrule(lr){2-7}
        & \textit{Avg. Score} & 0.71 & 0.72 & 0.72 & 0.80 & - \\
        \midrule

        % --- Level 4: Analysis  ---
        \multirow{4}{*}{\textbf{Analysis}} 
        & Semantic Tree Similarity         & 0.39 & 0.55 & 0.53 & 0.56 & - \\
        & Shape Consistency         & 0.54 & 0.74 & 0.72 & 0.75 & - \\
        & Structural Clarity Score               & 0.84 & 0.75 & 0.84 & 0.79 & 0.89 \\
        \cmidrule(lr){2-7}
        & \textit{Avg. Score} & 0.59 & 0.68 & 0.70 & 0.70 & - \\
        \midrule

        % --- Level 5: Evaluation ---
        \multirow{2}{*}{\textbf{Evaluation}} 
        & \makecell[l]{Critical Analysis Alignment (F1)}          & 0.13 & 0.08 & 0.11 & 0.13 & - \\
        \cmidrule(lr){2-7}
        & \textit{Avg. Score} & 0.13 & 0.08 & 0.11 & 0.13 & - \\
        \midrule

        % --- Level 6: Creation  ---
        \multirow{3}{*}{\textbf{Creation}} 
        & Framework Novelty              & 0.44 & 0.41 & 0.53 & 0.41 & - \\
        & \makecell[l]{Research Outlook Quality}               & 0.51 & 0.47 & 0.56 & 0.48 & - \\
        \cmidrule(lr){2-7}
        & \textit{Avg. Score} & 0.48 & 0.44 & 0.55 & 0.45 & - \\
        \bottomrule
    \end{tabular}
\end{table*}


Our comprehensive evaluation, with detailed results in Table~\ref{tab:main_tab} and a high-level summary in Figure~\ref{fig:all_levels}, reveals a substantial capability gap between four leading ASG systems and human experts. While LLM systems show competence in structured, formal tasks, they fall significantly short in higher-order cognitive dimensions requiring deep understanding and critical thinking.

\paragraph{Lower-Order Skills (Memory \& Comprehension).}
At the foundational levels, all systems exhibit significant shortcomings in factuality and faithfulness, with average Memory scores below 0.2. A severe inability to recall core, high-impact literature is evident, as no system's High-Impact Reference Coverage (HIRC) F1 score exceeded 0.07. Low Factual Consistency scores further expose issues with hallucination. While Entity Fidelity F1 scores were low, the high DS scores (0.69--0.81) suggest that when entities are correctly identified, their discussion frequency is human-like. Performance in Comprehension was slightly better (average scores 0.43--0.57), with models generally able to identify main topics (high OTC scores). However, universally low Citation Faithfulness (CF) scores highlight a critical failure to accurately understand and paraphrase source material.


\paragraph{Mid-Level Skills (Application \& Analysis).}
In the mid-level tiers, ASG systems reveal a stark divide between superficial formatting and deep structural understanding, demonstrating a tendency for  ``form over substance''. The Application tier is particularly illustrative of this split. While systems achieve near-perfect scores in low-level procedural tasks like Formatting Integrity (FMI) and Document Structure Integrity (DSI), their ability to apply a coherent high-level framework is notably weak. With the exception of LLMxMapReduce-V2, all systems scored well below the 0.5 human-parity baseline on Framework Application (FAP), indicating they are excellent ``format organizers'' but poor ``architectural thinkers''. In Analysis, high Structural Clarity (SCS) scores confirm that the generated outlines are internally coherent. However, this contrasts with moderate Semantic Tree Similarity (STS) scores (0.39--0.56), indicating that while structurally sound, the systems' knowledge organization still differs significantly from an expert's.

\paragraph{Higher-Order Skills (Evaluation \& Creation).}
The most critical bottlenecks appear at the highest cognitive levels. In Evaluation, all systems performed exceptionally poorly, with average scores clustering around 0.1, providing strong evidence that they almost completely lack human-aligned critical judgment. The gap is even more pronounced in Creation. Contrary to being creative, most systems failed to reach the 0.5 human-expert parity baseline, with average scores as low as 0.44. This demonstrates a significant and quantifiable lack of creative ability, rather than an ``illusion of creativity''. Only LLMxMapReduce-V2 managed to slightly exceed this baseline, suggesting a nascent but still limited creative capacity.

\paragraph{Overall System Comparison.}
No single system dominates across all dimensions. SurveyX emerges as the strongest all-around system for processing and structuring information, achieving the top scores in Application (0.80) and Comprehension (0.57), and tying for the lead in Analysis (0.70). In contrast, LLMxMapReduce-V2 displays a more specialized strength: it is the only system to surpass the human-parity baseline in both Framework Application (0.52) and Creation (0.55), but its performance in lower-order skills like Memory is weaker. SurveyForge demonstrates solid mid-range performance, particularly in Comprehension and Analysis, but does not lead in any single category. Overall, our results profile current ASG systems as excellent ``structural mimics'' but deficient ``knowledge integrators'' and genuine ``idea creators''. A significant journey remains before they can be considered reliable intellectual partners.

\paragraph{Efficiency Analysis.} We also examined operational efficiency (see Appendix Table~\ref{tab:appendix_supplementary}). The results highlight distinct resource profiles: LLMxMapReduce-V2 is the most token-intensive ($>$72k tokens), whereas SurveyX exhibits the highest latency (126 min). In comparison, SurveyForge demonstrates balanced efficiency, aligning with human-level length ($\approx$26k tokens) while accelerating the writing process by nearly $40\times$.

\section{Conclusion and Outlook}
We introduced Bloom-Eval, a hierarchical framework that reveals a critical gap in automatic survey generation (ASG): current systems are proficient ``format organizers'' but fail at higher-order cognitive tasks like critical analysis and creativity. This finding suggests that without a change in evaluation, the field risks optimizing for well-structured yet intellectually shallow content. Our work calls for a research paradigm shift from ``formal mimicry'' to ``cognitive deepening''. Bloom-Eval provides the diagnostic compass to guide this journey toward creating truly intelligent systems.

\section*{Limitations}
While Bloom-Eval offers a more systematic and fine-grained framework for evaluating ASG systems, we acknowledge several limitations in our framework design and experimental methodology. First, regarding the evaluation framework itself, although our proposed GRADE approach is designed to enhance transparency, the assessment of abstract, high-order abilities such as ``creativity'' ultimately relies on the judgment of a Large Language Model, which cannot fully escape its inherent subjective preferences. Second, our evaluation is constrained by the scope of the constructed benchmark. Despite our efforts to ensure its scale and interdisciplinary nature, the existing corpus of over 3,000 papers and 20 evaluation topics cannot fully represent the entire spectrum of scientific research domains.

% \section*{Acknowledgments}

% This document has been adapted
% by Steven Bethard, Ryan Cotterell and Rui Yan
% from the instructions for earlier ACL and NAACL proceedings, including those for
% ACL 2019 by Douwe Kiela and Ivan Vuli\'{c},
% NAACL 2019 by Stephanie Lukin and Alla Roskovskaya,
% ACL 2018 by Shay Cohen, Kevin Gimpel, and Wei Lu,
% NAACL 2018 by Margaret Mitchell and Stephanie Lukin,
% Bib\TeX{} suggestions for (NA)ACL 2017/2018 from Jason Eisner,
% ACL 2017 by Dan Gildea and Min-Yen Kan,
% NAACL 2017 by Margaret Mitchell,
% ACL 2012 by Maggie Li and Michael White,
% ACL 2010 by Jing-Shin Chang and Philipp Koehn,
% ACL 2008 by Johanna D. Moore, Simone Teufel, James Allan, and Sadaoki Furui,
% ACL 2005 by Hwee Tou Ng and Kemal Oflazer,
% ACL 2002 by Eugene Charniak and Dekang Lin,
% and earlier ACL and EACL formats written by several people, including
% John Chen, Henry S. Thompson and Donald Walker.
% Additional elements were taken from the formatting instructions of the \emph{International Joint Conference on Artificial Intelligence} and the \emph{Conference on Computer Vision and Pattern Recognition}.

% Bibliography entries for the entire Anthology, followed by custom entries
%\bibliography{anthology,custom}
% Custom bibliography entries only
\bibliography{custom}

\appendix


\section{Implementation Details of Evaluation Methodologies}
\label{sec:detail_Evaluation_Methodologies}

This appendix presents the mathematical formulations and implementation details of the evaluation methodologies introduced in the Bloom-Eval benchmark.

\subsection{Distributional Similarity (DS)}
\label{sec:ds_metric}

To quantify the distributional alignment between sets (e.g., entities or topics), we propose a composite metric named Distributional Similarity (DS). We model item frequencies as two probability distributions, $P$ and $Q$. Frequency counts from the expert ($C_\text{E}$) and generated ($C_\text{G}$) documents are normalized into the probability mass functions $P$ and $Q$ as follows:
{
\setlength{\abovedisplayskip}{6pt}%
\setlength{\belowdisplayskip}{6pt}%
\begin{equation}
p_i = \frac{c_{\text{E},i}}{\sum_{j=1}^{n} c_{\text{E},j}} \quad \text{and} \quad q_i = \frac{c_{\text{G},i}}{\sum_{j=1}^{n} c_{\text{G},j}}.
\end{equation}
}
The DS score is then calculated as a multi-metric fusion of three distances empirically validated in NLP contexts~\cite{balandranocoronel2016using, lebret2014word, ji2023tailoring}: the Jensen-Shannon Divergence $d_{\mathrm{JS}}$, the Hellinger Distance $d_{\mathrm{H}}$, and the Total Variation Distance $d_{\mathrm{TV}}$. Each distance is converted to a similarity score between $0$ and $1$, and then averaged. The Hellinger and Total Variation distances are used directly as their values naturally fall within the $[0, 1]$ range. The Jensen-Shannon Divergence is first normalized by its theoretical maximum:
{
\setlength{\abovedisplayskip}{6pt}%
\setlength{\belowdisplayskip}{6pt}%
\begin{equation}
\hat{d}_{\mathrm{JS}} = \frac{d_{\mathrm{JS}}(P, Q)}{\ln 2}.
\end{equation}
}
The final DS score is then computed as
{
\setlength{\abovedisplayskip}{6pt}%
\setlength{\belowdisplayskip}{6pt}%
\begin{equation*}
\mathrm{DS}(P, Q) = \frac{1}{3} \left( (1 - \hat{d}_{\mathrm{JS}}) + (1 - d_{\mathrm{H}}) + (1 - d_{\mathrm{TV}}) \right).
\end{equation*}
}
A higher DS score indicates greater similarity, with 1 signifying identical distributions.


\subsection{Generative Rubric Adaptive Differential Evaluation (GRADE)}
\label{sec:grade}

\begin{enumerate}[leftmargin=0.5cm, itemsep=0pt, parsep=0pt]
    \item \textbf{Rubric Generation:} A ``Judge LLM'' first deconstructs the evaluation task by generating a specific rubric. This yields a set of $k$ criteria $C=\{c_1, \dots, c_k\}$ and corresponding weights $W=\{w_1, \dots, w_k\}$, where $\sum w_i = 1$. This step externalizes the LLM's reasoning, making it explicit and auditable.

    \item \textbf{Differential Scoring:} The LLM then uses this public rubric to perform a comparative evaluation, assigning a score $s_i \in [0, 10]$ to each criterion for both the generated survey and a human expert's reference.
\end{enumerate}

An intermediate score for any document is calculated as the weighted sum of its criterion scores:
{
\setlength{\abovedisplayskip}{6pt}%  % 缩小公式上方间距
\setlength{\belowdisplayskip}{6pt}%  % 缩小公式下方间距
\begin{equation}
\label{eq:grade_intermediate_score}
\text{Score}_{\text{doc}} = \sum_{i=1}^{k} s_i \cdot w_i.
\end{equation}
}
The final GRADE score is then computed as a relative measure between the generated intermediate score $\text{Score}_{\text{Generated}}$ and the human's score $\text{Score}_{\text{Human}}$:
{
\setlength{\abovedisplayskip}{6pt}%  % 缩小公式上方间距
\setlength{\belowdisplayskip}{6pt}%  % 缩小公式下方间距
\begin{equation}
\label{eq:grade_final_score}
\text{GRADE Score} = \frac{\text{Score}_{\text{Generated}}}{\text{Score}_{\text{Generated}} + \text{Score}_{\text{Human}}}.
\end{equation}
}
A score of 0.5 indicates parity with the human expert, while a score greater than 0.5 suggests superior performance by the ASG system on that specific abstract ability. This process transforms the LLM from an unaccountable ``black-box'' judge into a tool for structured, auditable assessment. 

% By anchoring the final score to an explicit set of criteria, GRADE significantly enhances the reliability of evaluating complex cognitive skills like Framework Application and Framework Novelty.

\section{Implementation Details of Hierarchical Metrics}
\label{sec:appendix_hm}

This section provides the specific mathematical formulations and implementation parameters for the diagnostic metrics across the six cognitive tiers.

\subsection{Level 1: Memory Metrics}

\paragraph{Entity Fidelity (EFid).}
To assess domain knowledge retention, we employ a Judge LLM to extract fine-grained entities from both the generated and expert surveys, specifically targeting four categories: methods, models, datasets, and evaluation metrics. Based on these extractions, we calculate the standard F1-score to quantify the set overlap between the generated and expert entities, as well as the Distributional Similarity (DS) score (defined in Appendix~\ref{sec:ds_metric}) to evaluate the alignment of frequency vectors for shared entities.

\paragraph{High-Impact Reference Coverage (HIRC).}
This metric measures the recall of seminal works. Citation counts for all references in the expert bibliography are retrieved via Google Scholar. We define the set of high-impact papers $R_{\text{high}}$ as references with citation counts $c > 50$. HIRC is calculated as the recall rate:
{
\setlength{\abovedisplayskip}{3pt}
\setlength{\belowdisplayskip}{3pt}
\begin{equation}
    \text{HIRC} = \frac{|R_{\text{gen}} \cap R_{\text{high}}|}{|R_{\text{high}}|},
\end{equation}
}
where $R_{\text{gen}}$ is the set of references cited in the generated survey. This metric assesses whether the generated content preserves the cornerstone research that constitutes the domain's intellectual heritage.

\paragraph{Factual Consistency (FCons).} To distinguish faithful recall from hallucination, we implement a two-stage LLM-based verification pipeline that begins by extracting atomic factual claims from the generated text. A Judge LLM then evaluates each claim through semantic matching against the expert survey to determine entailment, verifying support from the ground truth. 


\subsection{Level 2: Comprehension Metrics}
\label{sec:app_l2}

\paragraph{Outline Topic Coverage (OTC).}
We employ an LLM to extract all headings from the outlines of both the generated and expert surveys. Subsequently, the LLM performs semantic matching between the two sets of headings to identify overlapping topics. Based on the matched set, we compute the F1-score.

\paragraph{Citation Faithfulness (CF).}
We adopt the Citation Recall and Precision metrics from~\cite{wang2024autosurvey} to distinguish faithful interpretation from hallucination. Specifically, for each citation, we retrieve the title and abstract of the referenced paper. A Judge LLM then determines whether the specific statement in the survey is textually supported by the retrieved reference context. Based on these judgments, we compute the final F1-score.

\paragraph{Thematic Focus Similarity (TFSim).}
TFSim measures whether the model captures the same research clusters as experts. We utilize the nomic-embed-text-v1 model~\cite{nussbaum2025nomic} to generate semantic embeddings for the concatenated titles and abstracts of all cited references in both bibliographies. These embeddings are then clustered to identify shared research themes. Based on the clustering results, we calculate the F1-score for theme overlap and the Distributional Similarity (DS) score for the frequency distribution of these themes.

\paragraph{Thematic Balance (TBal).}
Using the topic counts derived from the bibliography clustering analysis, TBal quantifies the balance of topic coverage within a single survey. We apply the standard Gini Coefficient (G) \cite{gini1912variabilita} to capture the distributional inequality. The final metric is defined as \(1 - G\):
    {
    \setlength{\abovedisplayskip}{6pt}%  % 缩小公式上方间距
    \setlength{\belowdisplayskip}{6pt}%  % 缩小公式下方间距
    \begin{equation}
    G = \frac{\sum_{i=1}^{n} \sum_{j=1}^{n} |x_i - x_j|}{2n \sum_{i=1}^{n} x_i} ,
    \end{equation}
    }
    where $x_i$ is the number of references in theme $i$. A higher TBal score thus reflects a more comprehensive research focus.

\subsection{Level 3: Application Metrics}
\paragraph{Formatting Integrity (FMI).} 
    This metric assesses the consistency between in-text citations and the bibliography to detect errors such as undefined citation markers and extraneous bibliographic entries. It is calculated as the Jaccard similarity between the set of citation identifiers found in the text \(C_{\text{cited}}\) and those defined in the bibliography \(R_{\text{defined}}\):
    {
    \setlength{\abovedisplayskip}{6pt}%  % 缩小公式上方间距
    \setlength{\belowdisplayskip}{6pt}%  % 缩小公式下方间距
    \begin{equation}
    \text{FMI} = \frac{|C_{\text{cited}} \cap R_{\text{defined}}|}{|C_{\text{cited}} \cup R_{\text{defined}}|}.
    \end{equation}
    }

\paragraph{Document Structure Integrity (DSI).} This checklist-based metric measures the presence of pre-defined, essential sections $S_{\text{req}}$ (e.g., Abstract, Introduction, Conclusion, References) within the set of sections extracted from the generated document, denoted as $S_{\text{found}}$. It verifies whether the model correctly constructs the necessary skeleton of a formal survey paper. The score is calculated as:
    {
    \setlength{\abovedisplayskip}{6pt}%  % 缩小公式上方间距
    \setlength{\belowdisplayskip}{6pt}%  % 缩小公式下方间距
    \begin{equation}
        \text{DSI} = \frac{\sum_{s \in S_{\text{req}}} \mathbb{I}(s \in S_{\text{found}})}{|S_{\text{req}}|},
    \end{equation}
    }
    where \(\mathbb{I}(\cdot)\) is an indicator function.


\paragraph{Framework Application (FAP).}
To evaluate the application of organizing paradigms, FAP adopts the GRADE approach. The Judge LLM is first instructed to synthesize bespoke, topic-specific scoring rubrics regarding the selection and consistent implementation of recognized taxonomies (e.g., chronological, methodological). Subsequently, utilizing these generated rubrics, the LLM evaluates both the ASG-generated survey and the human expert survey to assess the transfer of abstract organizational logic, ultimately deriving the relative GRADE score.

\subsection{Level 4: Analysis Metrics}
\paragraph{Semantic Tree Similarity (STS).}
This metric quantifies the hierarchical and semantic similarity between the generated outline (Tree A) and the expert outline (Tree B). We parse both outlines into tree structures and compute their similarity using the Tree Edit Distance (TED) algorithm. The cost for node insertion and deletion is 1, while the update cost for a node (i.e., a section heading) is its semantic distance, defined as \(1 - \cos(\theta)\) based on vector embedding similarity. The final normalized score is calculated as:
{
\setlength{\abovedisplayskip}{6pt}%  % 缩小公式上方间距
\setlength{\belowdisplayskip}{6pt}%  % 缩小公式下方间距
\begin{equation}
\text{Similarity} = 1 - \frac{\text{TED}(A, B)}{|A| + |B|},
\end{equation}
}
where \(|A|\) and \(|B|\) are the number of nodes in each tree. A score approaching 1 indicates high structural and semantic congruence. By modeling outlines as semantic trees and computing their edit distance, this metric verifies whether the model organizes sub-topics with relational logic (e.g., parent-child relationships) consistent with human experts.


\paragraph{Shape Consistency (SCons).}
This diagnostic metric compares the structural style of two outlines by quantifying the similarity of their depth and breadth. It is calculated as the geometric mean of two ratios: \textbf{Depth Consistency (DC)} and \textbf{Breadth Consistency (BC)}. Let the depths of the two outlines be $d_1$ and $d_2$, and their total node counts be $n_1$ and $n_2$. The consistency ratios are defined as the smaller value divided by the larger value for each dimension:
{
\setlength{\abovedisplayskip}{6pt}%
\setlength{\belowdisplayskip}{6pt}%
\begin{equation}
    \text{DC} = \frac{\min(d_1, d_2)}{\max(d_1, d_2)} \quad \text{and} \quad \text{BC} = \frac{\min(n_1, n_2)}{\max(n_1, n_2)}.
\end{equation}
}
The final Shape Consistency score is then computed as:
{
\setlength{\abovedisplayskip}{6pt}%
\setlength{\belowdisplayskip}{6pt}%
\begin{equation}
    \text{SCons} = \sqrt{\text{DC} \times \text{BC}}.
\end{equation}
}


\paragraph{Structural Clarity Score (SCS).}
This metric assesses an outline's logical coherence by quantifying improper content duplication. We use an LLM to identify semantically equivalent heading pairs that do not share an immediate parent, a rule designed to prevent misidentifying normal sibling sections as redundant. These pairs are counted as redundant \(N_{\text{redundant}}\). The clarity score is then calculated as:
{
\setlength{\abovedisplayskip}{6pt}%  % 缩小公式上方间距
\setlength{\belowdisplayskip}{6pt}%  % 缩小公式下方间距
\begin{equation}
    \text{SCS} = 1 - \frac{N_{\text{redundant}}}{N_{\text{total}}},
\end{equation}
}
where \(N_{\text{total}}\) is the total number of sections. A score near 1 suggests a clear structure with minimal content overlap.

\subsection{Level 5: Evaluation Metrics}
\label{sec:app_l5}

\paragraph{Critical Analysis Alignment (CAA).}
This metric quantifies the extent to which the model's critical stance aligns with expert consensus. First, a Judge LLM is employed to extract explicit critical statements from both the generated and expert surveys. The prompt specifically instructs the model to identify sentences that describe limitations, methodological flaws, or comparative disadvantages of existing works. Second, to measure the overlap, a Judge LLM performs semantic matching between the two sets of extracted statements. For each critical statement in the generated set, the Judge determines if it is semantically equivalent to (or entailed by) any statement in the expert set.

Based on the number of LLM-verified matches, we report the F1-score to verify if the model's critical insights are both accurate (Precision) and exhaustive (Recall) relative to human experts. 

\subsection{Level 6: Creation Metrics}
\label{sec:app_l6}

\paragraph{Framework Novelty (FNov).}
To evaluate the framework novelty of the survey, FNov adopts the GRADE approach. The Judge LLM is first instructed to synthesize bespoke, topic-specific scoring rubrics regarding framework innovation. Subsequently, utilizing these generated rubrics, the LLM evaluates both the ASG-generated survey and the human expert survey, ultimately deriving the relative GRADE score.

\paragraph{Research Outlook Quality (ROQ).}
Similarly, ROQ utilizes the GRADE approach to assess the quality of future research directions. The Judge LLM is prompted to synthesize bespoke, topic-specific scoring rubrics concerning the foresight and value of the proposed research agenda. Based on these dynamic criteria, the LLM scores both the generated and expert surveys to calculate the final relative GRADE score.

\section{Experimental Topic Details}
To ground our evaluation in a rigorous and representative context, we selected 20 topics for our experiments. These topics correspond to the titles of the 20 most-cited survey papers published between 2023 and 2025 within our benchmark corpus, ensuring relevance to the current scientific frontier. The selection is evenly split, with 10 topics drawn from the Computer Science domain and 10 from a diverse range of General Science disciplines. This balanced approach allows for a robust assessment of the generalization capabilities of the evaluated ASG systems. The complete list of these 20 topics, along with their publication details and citation counts, is provided in Table~\ref{tab:experimental_topics}.

\begin{table*}[h!]
\centering
\caption{The 20 experimental topics used for evaluation, selected based on the most-cited survey papers from our benchmark corpus (2023-2025).}
\label{tab:experimental_topics}
\footnotesize
\begin{tabularx}{\textwidth}{>{\raggedright\arraybackslash}X l c r}
\toprule
\textbf{Title} & \textbf{Venue} & \textbf{Year} & \textbf{Citations} \\
\midrule
\multicolumn{4}{c}{\textit{Computer Science Domain}} \\
\midrule
A Survey on Vision Transformer & TPAMI & 2023 & 3915 \\
A Survey on In-context Learning & EMNLP & 2023 & 2379 \\
Diffusion Models in Vision: A Survey & TPAMI & 2023 & 2037 \\
Domain Generalization: A Survey & TPAMI & 2023 & 1720 \\
Transformers in Time Series: A Survey & IJCAI & 2023 & 1458 \\
Generalized Out-of-Distribution Detection: A Survey & IJCV & 2024 & 1377 \\
A Comprehensive Survey of Continual Learning: Theory, Method and Application & TPAMI & 2024 & 1313 \\
Class-Incremental Learning: Survey and Performance Evaluation on Image Classification & TPAMI & 2023 & 1038 \\
Transfer Learning in Deep Reinforcement Learning: A Survey & TPAMI & 2023 & 1031 \\
Multimodal Learning With Transformers: A Survey & TPAMI & 2023 & 990 \\
\midrule
\multicolumn{4}{c}{\textit{General Science Domain}} \\
\midrule
Structural Mechanisms of NLRP3 Inflammasome Assembly and Activation & Immunology & 2023 & 781 \\
The New Economics of Industrial Policy & Economics & 2023 & 567 \\
Random Quantum Circuits & Condensed Matter Physics & 2023 & 494 \\
Biological Impacts of Marine Heatwaves & Marine Science & 2023 & 468 \\
How to Run Surveys: A Guide to Creating Your Own, Identifying Variation and Revealing the Invisible & Economics & 2023 & 443 \\
A Review of Generalizability and Transportability & Statistics and Its Application & 2023 & 387 \\
Non-Hermitian Topological Phenomena: A Review & Condensed Matter Physics & 2023 & 358 \\
Hepcidin and Iron in Health and Disease & Medicine & 2023 & 340 \\
Generalized Symmetries in Condensed Matter & Condensed Matter Physics & 2023 & 335 \\
Everything You Wanted to Know about Deep Eutectic Solvents but Were Afraid to Be Told & Chemical and Biomolecular Engineering & 2023 & 258 \\
\bottomrule
\end{tabularx}
\end{table*}



\section{Detailed Experimental Results}
\label{sec:appendix_detailed_results}

This section presents the complete and disaggregated results of our evaluation. For each of the six cognitive tiers in the Bloom-Eval framework, we provide a detailed table showing the Precision (P), Recall (R), F1-Score (F1), and Distributional Similarity (DS) for every metric where applicable. Additionally, we report on several supplementary metrics, including generation speed and economic cost, to provide a more holistic view of system performance.

\subsection{Core Cognitive Performance Metrics}
For each of the evaluated cognitive tiers, Table~\ref{tab:appendix_level1}, Table~\ref{tab:appendix_level2}, and Table~\ref{tab:appendix_level5} show the detailed Precision (P), Recall (R), F1-Score (F1), and Distributional Similarity (DS) for every metric where applicable.
% --- Level 1 Table ---
\begin{table*}[h!]
    \centering
    \caption{Detailed results for Level 1: Memory. EFid scores are in P/R/F1/DS format. HIRC and FCons are in P/R/F1 format.}
    \label{tab:appendix_level1}
    \begin{tabular}{llcccc}
        \toprule
        \textbf{Metric} & \textbf{Values} & \textbf{AutoSurvey} & \textbf{SurveyForge} & \textbf{LLMxMapReduce-V2} & \textbf{SurveyX} \\
        \midrule
        \multirow{4}{*}{EFid} 
            & Precision & 0.11 & 0.17 & 0.11 & 0.13 \\
            & Recall    & 0.10 & 0.11 & 0.17 & 0.12 \\
            & F1-Score  & 0.10 & 0.13 & 0.12 & 0.11 \\
            & DS-Score  & 0.81 & 0.73 & 0.69 & 0.78 \\
        \midrule
        \multirow{3}{*}{HIRC} 
            & Precision & 0.03 & 0.11 & 0.04 & 0.06 \\
            & Recall    & 0.05 & 0.06 & 0.01 & 0.04 \\
            & F1-Score  & 0.03 & 0.07 & 0.02 & 0.04 \\
        \midrule
        \multirow{3}{*}{FCons} 
            & Precision & 0.07 & 0.06 & 0.19 & 0.10 \\
            & Recall    & 0.10 & 0.09 & 0.07 & 0.06 \\
            & F1-Score  & 0.07 & 0.06 & 0.09 & 0.07 \\
        \bottomrule
    \end{tabular}
\end{table*}


% --- Level 2 Table ---
\begin{table*}[h!]
    \centering
    \caption{Detailed results for Level 2: Comprehension. TFSim is in P/R/F1/DS format. OTC and CF are in P/R/F1 format.}
    \label{tab:appendix_level2}
    \begin{tabular}{llcccc}
        \toprule
        \textbf{Metric} & \textbf{Values} & \textbf{AutoSurvey} & \textbf{SurveyForge} & \textbf{LLMxMapReduce-V2} & \textbf{SurveyX} \\
        \midrule
        \multirow{3}{*}{OTC} 
            & Precision & 0.29 & 0.52 & 0.41 & 0.43 \\
            & Recall    & 0.76 & 0.78 & 0.82 & 0.81 \\
            & F1-Score  & 0.40 & 0.61 & 0.53 & 0.56 \\
        \midrule
        \multirow{3}{*}{CF} 
            & Precision & 0.03 & 0.27 & 0.39 & 0.44 \\
            & Recall    & 0.03 & 0.26 & 0.23 & 0.32 \\
            & F1-Score  & 0.03 & 0.27 & 0.29 & 0.37 \\
        \midrule
        \multirow{4}{*}{TFSim} 
            & Precision & 0.35 & 0.68 & 0.73 & 0.62 \\
            & Recall    & 0.76 & 0.56 & 0.32 & 0.55 \\
            & F1-Score  & 0.43 & 0.53 & 0.41 & 0.52 \\
            & DS-Score  & 0.73 & 0.78 & 0.67 & 0.81 \\
        \midrule
        TBal & 1-Gini & 0.72 & 0.70 & 0.68 & 0.68 \\
        \bottomrule
    \end{tabular}
\end{table*}

% --- Level 5 Table ---
\begin{table*}[h!]
    \centering
    \caption{Detailed results for Level 5: Evaluation. Critical Analysis Alignment (CAA) scores are in P/R/F1 format.}
    \label{tab:appendix_level5}
    \begin{tabular}{llcccc}
        \toprule
        \textbf{Metric} & \textbf{Values} & \textbf{AutoSurvey} & \textbf{SurveyForge} & \textbf{LLMxMapReduce-V2} & \textbf{SurveyX} \\
        \midrule
        \multirow{3}{*}{CAA} 
            & Precision & 0.09 & 0.04 & 0.07 & 0.08 \\ 
            & Recall    & 0.35 & 0.38 & 0.45 & 0.43 \\ 
            & F1-Score  & 0.13 & 0.08 & 0.11 & 0.13 \\
        \bottomrule
    \end{tabular}
\end{table*}

\subsection{Supplementary Metrics: Efficiency and Readability}

To provide a more holistic view of system performance, we report on several supplementary metrics, with detailed results presented in Table~\ref{tab:appendix_supplementary}. These include the generation speed, total output token count, and a composite score for Readability.

The Readability (R) score is designed to measure stylistic similarity. It is not an absolute measure of simplicity but rather a score from 0 to 1 indicating how closely an ASG system's writing style mimics that of the human-written gold standard. By definition, the human-authored text has a perfect score of 1. The final score is a composite, calculated as the arithmetic mean of similarity scores from three established readability formulas: the Flesch-Kincaid Grade Level \cite{kincaid1975derivation}, the Gunning Fog index \citep{Gunning1952}, and the Coleman-Liau Index \citep{coleman1975computer}.

For each formula, a similarity ratio ($\text{sim}$) is first calculated by comparing the system's metric value ($v_{\text{generated}}$) to the human's ($v_{\text{human}}$):
\begin{align}
\text{sim}(v_{\text{human}}, v_{\text{generated}}) = \frac{\min(v_{\text{human}}, v_{\text{generated}})}{\max(v_{\text{human}}, v_{\text{generated}})}. \nonumber
\end{align}
This ensures that a larger divergence from the human baseline in either direction results in a lower score. The final Readability score ($R$) is the unweighted average of the three individual similarity scores ($\text{sim}_1, \text{sim}_2, \text{sim}_3$):
\begin{align}
R = \frac{1}{3} \sum_{i=1}^{3} \text{sim}_i.
\end{align}

This approach leverages these well-established formulas to quantify the stylistic similarity in terms of overall text complexity. Each formula computes a score, typically representing a U.S. grade level, by combining sentence-level features (e.g., average sentence length) with word-level features (e.g., syllable or character counts). By averaging the similarity scores derived from these distinct formulas, we mitigate the biases inherent in any single metric and obtain a more robust and comprehensive assessment of how closely the machine-generated text's complexity profile matches the human standard.

The evaluation of speed distinguishes between automated systems and human experts. For ASG systems, speed is measured as the time from the initial submission of a research topic to the final generation of the complete survey. For human experts, We follow the established estimation methodology \cite{wang2024autosurvey}.

% --- Supplementary Metrics Table ---
\begin{table*}[h!]
    \centering
    \caption{Supplementary metrics assessing efficiency and cost. Readability (R) is on a 0-1 scale. Speed (S) is in minutes. Tokens (T) is the total output count.}
    \label{tab:appendix_supplementary}
    \begin{tabular}{lccccc}
        \toprule
        \textbf{Metric} & \textbf{AutoSurvey} & \textbf{SurveyForge} & \textbf{LLMxMapReduce-V2} & \textbf{SurveyX} & \textbf{Human} \\
        \midrule
        Readability (R) & 0.73 & 0.71 & 0.75 & 0.68 & 1.00 \\
        Speed (S) & 12 & 11 & 60 & 126 & 434 \\
        Output Tokens (T) & 65,907 & 25,911 & 72,217 & 23,630 & 29,274 \\
        \bottomrule
    \end{tabular}
\end{table*}




\section{Comparative Performance Across Scientific Domains}
\label{sec:appendix_domain_comparison}

While many automatic survey generation (ASG) systems are capable of retrieving and processing information across multiple scientific disciplines, some systems, such as AutoSurvey \citep{wang2024autosurvey} and SurveyForge \citep{yan2025surveyforge}, are architecturally designed or configured to operate primarily within the Computer Science (CS) domain. This presents a critical question regarding their cross-domain generalization capabilities. To investigate this limitation, this section provides a detailed performance analysis of these two systems, contrasting their efficacy on their native CS topics with a broader set of General Science topics. The specific topics for the CS domain correspond to the first ten entries in our experimental corpus (see Table~\ref{tab:experimental_topics}), while the General Science topics are represented by the subsequent ten entries.

The following radar charts visualize the performance drop-off across key cognitive tiers of the Bloom-Eval framework. Figures \ref{fig:as_l1} to \ref{fig:as_grade} illustrate the domain-specific performance for AutoSurvey, while Figures \ref{fig:sf_l1} to \ref{fig:sf_grade} show the results for SurveyForge. This analysis quantifies the performance penalty of domain-specific architectures, highlighting the critical value of cross-disciplinary benchmarks for testing true generalization.

\begin{figure}[h!]
    \centering
    \includegraphics[width=0.4\textwidth]{latex/cs vs general/autosurvey_level1.png}
    \caption{The performance of AutoSurvey on Level 1 (Memory) metrics, comparing Computer Science and General Science domains.}
    \label{fig:as_l1}
\end{figure}

\begin{figure}[h!]
    \centering
    \includegraphics[width=0.4\textwidth]{latex/cs vs general/autosurvey_level2.png}
    \caption{The performance of AutoSurvey on Level 2 (Comprehension) metrics, comparing CS and General Science domains.}
    \label{fig:as_l2}
\end{figure}

\begin{figure}[h!]
    \centering
    \includegraphics[width=0.4\textwidth]{latex/cs vs general/autosurvey_level4.png}
    \caption{The performance of AutoSurvey on Level 4 (Analysis) metrics, comparing CS and General Science domains.}
    \label{fig:as_l4}
\end{figure}

\begin{figure}[h!]
    \centering
    \includegraphics[width=0.4\textwidth]{latex/cs vs general/autosurvey_grade_metrics.png}
    \caption{The performance of AutoSurvey on the GRADE-based metrics, comparing CS and General Science domains.}
    \label{fig:as_grade}
\end{figure}

The performance of AutoSurvey demonstrates a strong dependency on its native Computer Science domain, with a notable degradation observed when applied to General Science topics. This performance gap is most pronounced in knowledge-intensive, memory-based tasks (Level 1). Specifically, the system's ability to identify high-impact references (HIRC) is almost entirely nullified in the General Science domain, accompanied by significant reductions in entity fidelity and factual consistency. Similarly, comprehension-level metrics such as Citation Faithfulness (Level 2) also exhibit a marked decline. 

In stark contrast, the system's analytical capabilities for structuring content (Level 4) appear to be largely domain-agnostic. Metrics such as Semantic Tree Similarity, Shape Consistency, and Structural Clarity Score remain remarkably stable across both domains. For higher-order GRADE-based tasks, while AutoSurvey's ability to generate novel frameworks and high-quality research outlooks is strong within CS, this capacity is diminished in the broader scientific context. This suggests that the primary limitation stems not from the model's inherent reasoning or structural generation abilities, but rather from its retrieval mechanism, which is heavily restricted to the Computer Science literature.


\begin{figure}[h!]
    \centering
    \includegraphics[width=0.4\textwidth]{latex/cs vs general/surveyforge_level1.png}
    \caption{The performance of SurveyForge on Level 1 (Memory) metrics, comparing Computer Science and General Science domains.}
    \label{fig:sf_l1}
\end{figure}

\begin{figure}[h!]
    \centering
    \includegraphics[width=0.4\textwidth]{latex/cs vs general/surveyforge_level2.png}
    \caption{The performance of SurveyForge on Level 2 (Comprehension) metrics, comparing CS and General Science domains.}
    \label{fig:sf_l2}
\end{figure}

\begin{figure}[h!]
    \centering
    \includegraphics[width=0.4\textwidth]{latex/cs vs general/surveyforge_level4.png}
    \caption{The performance of SurveyForge on Level 4 (Analysis) metrics, comparing CS and General Science domains.}
    \label{fig:sf_l4}
\end{figure}

\begin{figure}[h!]
    \centering
    \includegraphics[width=0.4\textwidth]{latex/cs vs general/surveyforge_grade_metrics.png}
    \caption{The performance of SurveyForge on the GRADE-based metrics, comparing CS and General Science domains.}
    \label{fig:sf_grade}
\end{figure}

SurveyForge also exhibits a significant performance dependency on the Computer Science (CS) domain. The system's most critical failure outside of CS lies in its knowledge grounding and citation abilities. In the General Science domain, the capacity to identify high-impact references is severely diminished, and Citation Faithfulness drops by more than half, indicating a fundamental breakdown in connecting claims to correct, relevant sources. 

% SurveyForge exhibits performance degradation across multiple cognitive tiers when applied outside its native CS domain. This decline is evident in Memory (Level 1), Analysis (Level 4), the GRADE-based metrics, and particularly in Citation Faithfulness (Level 2). This widespread degradation strongly suggests that a retrieval mechanism optimized for Computer Science is the primary bottleneck, hindering the system's capabilities in cross-domain applications.

\section{Diagnostic Analysis of Performance Bottlenecks}
\label{sec:appendix_root_cause}

To better understand the capability gap identified in our experiments, specifically why ASG systems fail at higher-order cognitive tiers, we analyze the systemic bottlenecks of current architectures. 

\begin{itemize}[leftmargin=*, parsep=8pt]
    \item \textbf{Retrieval Misalignment.} Existing systems universally failed the High-Impact Reference Coverage (HIRC) metric, even in their native CS domain. Our analysis reveals that current retrieval algorithms (e.g., in AutoSurvey~\cite{wang2024autosurvey}) rely heavily on vector embedding similarity while ignoring bibliometric impact. This leads to the retrieval of ``semantically overlapping but mediocre'' papers, missing foundational works with high citation counts. Furthermore, the reliance on static, closed-source CS databases leads to a severe degradation in cross-disciplinary tasks. Future work must integrate bibliometric indicators (e.g., citation velocity) into the retrieval ranker. Our contribution of a cross-disciplinary benchmark is specifically designed to guide the field away from closed CS datasets toward robust, open-domain mechanisms.

    \item \textbf{Information Processing via ``Lossy Compression''.} Systems scored poorly in Creation and Evaluation tiers because they lack deep insight. To fit limited context windows, baselines employ aggressive compression strategies. For instance, SurveyX~\cite{liang2025surveyx} compresses papers into ``AttributeTrees'', and LLMxMapReduce-V2~\cite{wang2025llmmapreducev2} relies on ``Digests''. This processing strips away subtle derivations, negative results, and complex arguments, effectively depriving the LLM of the raw materials needed for deep synthesis and innovation. The field must move from lossy summarization to long-context architectures that can process high-fidelity full text, preserving the nuance required for higher-order cognition.

    \item \textbf{Underutilization of Meta-Knowledge.} Existing methods like SurveyForge~\cite{yan2025surveyforge} utilize related survey papers primarily to mimic their ``outline structure'', rather than understanding the intellectual lineage or core debates within the field. This strategy results in models that produce a ``flat'' listing of literature, failing to synthesize knowledge depth. 
\end{itemize}

\section{Reliability and Validation of Evaluation Metrics}
\label{sec:appendix_validation}

In this section, we provide a comprehensive validation of the Bloom-Eval framework through two complementary analyses: a global correlation study covering all cognitive tiers (Section~\ref{sec:human_correlation}), and a specific sensitivity analysis for the abstract GRADE metrics (Section~\ref{sec:grade_sensitivity}).

\subsection{Alignment with Human Judgments}
\label{sec:human_correlation}

To validate the overall effectiveness of Bloom-Eval, we assessed whether our automated metrics, ranging from deterministic algorithms to LLM-based evaluations, align with human consensus across the six cognitive tiers.

We recruited five evaluators to provide human judgments. The panel comprised three Ph.D. candidates and two Master's students. To ensure domain diversity, we selected the three most-cited topics from both the Computer Science and General Science domains (six topics in total). For each topic, we collected survey papers generated by four different ASG systems plus the expert-written references, resulting in a total of 30 papers. To ensure the human evaluation was both rigorous and tractable, we selected one representative metric for each cognitive tier that captures the core definition of that level for human assessment.

We calculated the Pearson and Spearman correlations between the scores of these selected representative metrics (computed using GPT-5-mini for GRADE) and the human ratings. As shown in Table~\ref{tab:human_corr_validation}, Bloom-Eval demonstrates strong alignment ($>0.7$) with human experts across all dimensions. This result confirms two key findings:
\begin{enumerate} 
\item \textbf{Validation of Engineered Metrics:} Our metrics (e.g., Factual Consistency, Citation Faithfulness) successfully capture the nuances of human assessment.

\item \textbf{Evaluator Reliability:} The choice of GPT-5-mini as the underlying model for GRADE provides a rigorous standard consistent with expert reviewers, justifying its use as the primary judge in our main experiments.
\end{enumerate}

\begin{table*}[t] 
    \centering
    \caption{Correlations between Bloom-Eval metrics and human expert ratings on a subset of 30 papers.}
    \label{tab:human_corr_validation}
    \small
    \setlength{\tabcolsep}{10pt} % 既然跨栏了，可以让间距宽一点，更好看
    \begin{tabular}{llcc}
        \toprule
        Cognitive Tier & Automated Metric & Pearson & Spearman \\
        \midrule
        Level 1: Memory & Factual Consistency & 0.76 & 0.70 \\
        Level 2: Comprehension & Citation Faithfulness & 0.92 & 0.92 \\
        Level 3: Application & Framework Application (GRADE) & 0.83 & 0.84 \\
        Level 4: Analysis & Structural Clarity Score & 0.75 & 0.76 \\
        Level 5: Evaluation & Critical Analysis Alignment & 0.79 & 0.84 \\
        Level 6: Creation & Research Outlook Quality (GRADE) & 0.70 & 0.75 \\
        \bottomrule
    \end{tabular}
\end{table*}

\subsection{Sensitivity Analysis of the GRADE Approach}
\label{sec:grade_sensitivity}

Notwithstanding the general consistency demonstrated above, the inherent subjectivity of abstract reasoning metrics (specifically in Levels 3 and 6) warrants deeper scrutiny regarding their sensitivity to the choice of the Judge LLM.

To investigate this, we conducted a sensitivity analysis on the Framework Application (FAP) metric, a key indicator of a system's ability to apply knowledge.

In this experiment, we selected three distinct large language models, GPT-4.1-mini, GPT-5-mini, and GPT-4o-mini, to serve as independent judges. Each model evaluated the same 20 generated surveys against their human-expert counterparts and produced a GRADE score for the FAP metric. The results reveal a crucial finding: the evaluation outcome is sensitive to the choice of the Judge LLM, exposing significant ``Rater Bias'' inherent in the ``LLM-as-a-Judge'' paradigm. This finding does not weaken our framework. On the contrary, it highlights the critical necessity of our transparent, rubric-based design.

\subsubsection{Divergence in Scoring Bias and Qualitative Conclusions}

The most immediate finding is the profound difference in the scoring tendencies of the three models, as illustrated in Figure~\ref{fig:grade_sensitivity_boxplot}.

\begin{figure}[h!]
    \centering
    \includegraphics[width=0.9\columnwidth]{latex/related/related_1.png} 
    \caption{Distribution of Framework Application (FAP) scores across 20 tasks, as judged by three different LLMs. The red dashed line indicates the human parity baseline. The plot reveals starkly different scoring biases: optimistic (GPT-4.1-mini), pessimistic (GPT-5-mini), and unstable (GPT-4o-mini).}
    \label{fig:grade_sensitivity_boxplot}
\end{figure}

The analysis reveals three distinct rater profiles:
\begin{itemize}[leftmargin=*,itemsep=2pt,topsep=3pt]
    \item \textbf{The Optimistic Rater (GPT-4.1-mini):} This model exhibits a strong positive bias, with a mean score of 0.554. In 18 out of 20 cases (90\%), it judged the generated survey's framework application as superior to the human expert's. Its scoring is also relatively consistent, with a small standard deviation (0.042).
    
    \item \textbf{The Pessimistic Rater (GPT-5-mini):} In stark contrast, this model shows a consistent negative bias. Its mean score is only 0.407, and it judged the system's performance to be inferior to the human baseline in all 20 cases (100\%). It is the most consistent rater, with the smallest standard deviation (0.035).
    
    \item \textbf{The Unstable Rater (GPT-4o-mini):} This model lacks a clear, consistent standard. Its mean score of 0.506 hovers around the human parity line. Critically, its standard deviation (0.089) is more than double that of the other two models, indicating highly variable and unpredictable judgments.
\end{itemize}
This divergence is alarming: a researcher's conclusion about a system's capability on this metric is predetermined by their choice of Judge LLM.

\paragraph{Directional Agreement.} We calculated the rate at which any two models agreed on whether a score was $\ge 0.5$ or $< 0.5$. The results in Table~\ref{tab:directional_agreement} are presented below.

\begin{table}[h!]
\centering
\caption{Directional Agreement Rate between Judge LLMs.}
\label{tab:directional_agreement}
\begin{tabular}{@{}lc@{}}
\toprule
\textbf{Model Pair} & \textbf{Agreement Rate} \\
\midrule
GPT-4.1-mini vs. GPT-5-mini & 10\% \\
GPT-4.1-mini vs. GPT-4o-mini & 60\% \\
GPT-5-mini vs. GPT-4o-mini & 50\% \\
\bottomrule
\end{tabular}
\end{table}
The minimal 10\% agreement rate between the optimistic and pessimistic raters signifies a fundamental breakdown in consensus, highlighting a systematic opposition in their evaluative judgments. Furthermore, the moderate agreement levels observed in the other pairings remain below the threshold required for reliable inter-rater consistency.


\paragraph{Analysis of Correlational Agreement.}
To assess a deeper level of inter-rater reliability, we performed a correlational analysis. We investigated two forms of correlation: linear association using Pearson's correlation and ordinal consistency using Spearman's rank-order correlation.

\begin{figure}[h!]
    \centering
    \includegraphics[width=1\columnwidth]{latex/related/related_3.png} %<-- USER: Use the Pearson heatmap image file
    \caption{Pearson correlation heatmap of FAP scores.}
    \label{fig:pearson_correlation_heatmap}
\end{figure}

First, we examined the linear relationship between the models' scores. As shown in Figure~\ref{fig:pearson_correlation_heatmap}, the scores exhibit a weak to moderate positive correlation, with coefficients ranging from 0.47 to 0.67. While this suggests a general tendency for scores to move in the same direction, it does not sufficiently establish a robust, shared evaluative standard.

\begin{figure}[h!]
    \centering
    \includegraphics[width=1\columnwidth]{latex/related/related_4.png} %<-- USER: Use the Spearman heatmap image file
    \caption{Spearman's rank-order correlation heatmap.}
    \label{fig:spearman_correlation_heatmap}
\end{figure}

We also assessed the consistency of ordinal judgments using Spearman's rank correlation. The results, visualized in Figure~\ref{fig:spearman_correlation_heatmap}, reveal a significant lack of consensus on relative quality. The correlation coefficients are uniformly low, with the value between the optimistic rater (GPT-4.1-mini) and the pessimistic rater (GPT-5-mini) being a mere 0.35. This weak rank-order correlation is a key finding: it indicates that the models do not share a stable, underlying logic for what constitutes a ``better'' or ``worse'' application of a framework. Their disagreement is therefore not just a matter of scale, but a fundamental divergence in evaluative criteria.

\paragraph{The Value of Transparency in a Biased World.}

The severe sensitivity to rater selection presents a major challenge to the ``LLM-as-a-Judge'' paradigm. However, rather than invalidating our framework, this finding powerfully underscores the value of the GRADE approach's core design principle: transparency.

\begin{enumerate}
    \item \textbf{Rubrics as an Auditable Artifact:} Since any Judge LLM is inherently biased, we cannot trust a single, opaque score. The first step of GRADE forces the LLM to generate an explicit, weighted rubric. This rubric becomes a public, auditable artifact. It shifts the debate from an inscrutable final score to a transparent set of criteria that can be debated, refined, and analyzed, thus mitigating the black-box problem.

    \item \textbf{A Methodological Guideline:} Our results serve as a critical methodological warning for the field. Any study employing an LLM-as-a-Judge must report the specific model version used and acknowledge that results may not be directly comparable to studies using different judges. Moreover, they should report quantitative evidence of the judge's alignment with human experts (e.g., correlation coefficients). 
    
    \item \textbf{Ensuring Internal Consistency:} Given the strong consistency between GPT-5-mini and human expert judgments, we ultimately selected GPT-5-mini as the judge.
    
\end{enumerate}

This sensitivity analysis confirms that no LLM is an objective, ``god-like'' evaluator. Recognizing this reality, the GRADE approach provides the necessary methodological scaffolding through the transparency of explicit rubrics, making LLM-based evaluation more rigorous, interpretable, and scientifically sound.

\subsubsection{Case Study: Deconstructing Rater Bias Through Transparent Rubrics}
\label{sec:case_study_rater_bias}

To understand the origin of rater bias, we conducted a case study on a task where the Judge LLMs profoundly disagreed. Our analysis of their explicit, model-generated rubrics reveals that the scoring divergence stems from three distinct rater archetypes with different ``academic values'': GPT-4o-mini, the ``Pragmatic Balancer'', seeks a compromise between structure and content (Table~\ref{tab:rubric_gpt4o}); GPT-5-mini, the ``Rigorous Methodologist'', demands scientific rigor like a peer reviewer (Table~\ref{tab:rubric_gpt5}); and GPT-4.1-mini, the ``Structuralist'', prioritizes formal organization like an editor (Table~\ref{tab:rubric_gpt41}). This vividly demonstrates that rater bias is not random error but a systematic manifestation of different latent definitions of ``quality''. Crucially, this analysis is only possible due to our transparent, rubric-based methodology, which transforms the LLM's ``black box'' reasoning into a ``glass box'' and makes the evaluation process itself an object of scientific inquiry.

% 这是一个统一的、包含所有内容的父表格环境
\begin{table*}[h]
\centering
\caption{Evaluation rubrics generated by the three different Judge LLMs.}
\label{tab:all_rubrics}
\footnotesize

% --- 子表格 a: GPT-4o-mini ---
\begin{subtable}{\textwidth}
    \centering
    % 这是子表格的标题，它将被编号为 "(a)"
    \caption{Evaluation rubric generated by GPT-4o-mini.}
    \label{tab:rubric_gpt4o}
    \begin{tabularx}{0.8\textwidth}{X c}
        \toprule
        \textbf{Criterion} & \textbf{Weight} \\
        \midrule
        Relevance and Appropriateness of Framework & 30\% \\
        Clarity and Definition of Framework Components & 20\% \\
        Consistent Application of Framework & 20\% \\
        Analytical Depth and Insight & 20\% \\
        Domain-Specific Adaptability & 10\% \\
        \bottomrule
    \end{tabularx}
\end{subtable}

\vspace{1em} % 增加一点垂直间距

% --- 子表格 b: GPT-5-mini ---
\begin{subtable}{\textwidth}
    \centering
    % 这是子表格的标题，它将被编号为 "(b)"
    \caption{Evaluation rubric generated by GPT-5-mini.}
    \label{tab:rubric_gpt5}
    \begin{tabularx}{0.8\textwidth}{X c}
        \toprule
        \textbf{Criterion} & \textbf{Weight} \\
        \midrule
        Clarity and Relevance of the Chosen Framework & 22\% \\
        Taxonomy Coverage and Granularity & 18\% \\
        Consistent Application of the Framework & 17\% \\
        Analytical Depth and Trade-off Analysis Enabled by the Framework & 16\% \\
        Mapping to Tasks, Datasets, and Evaluation Protocols & 10\% \\
        Evidence-based Classification and Traceability & 9\% \\
        Actionability: Gaps, Best Practices, and Research Agenda & 8\% \\
        \bottomrule
    \end{tabularx}
\end{subtable}

\vspace{1em} % 增加一点垂直间距

% --- 子表格 c: GPT-4.1-mini ---
\begin{subtable}{\textwidth}
    \centering
    % 这是子表格的标题，它将被编号为 "(c)"
    \caption{Evaluation rubric generated by GPT-4.1-mini.}
    \label{tab:rubric_gpt41}
    \begin{tabularx}{0.8\textwidth}{X c}
        \toprule
        \textbf{Criterion} & \textbf{Weight} \\
        \midrule
        Clarity and Relevance of Framework & 40\% \\
        Consistent Application of Framework & 35\% \\
        Analytical Depth Enabled by Framework & 25\% \\
        \bottomrule
    \end{tabularx}
\end{subtable}

\end{table*}

\subsubsection{Case Study: Auditability Comparison}
\label{appendix:g22_comparison}

To demonstrate how the GRADE approach transforms a ``black-box'' evaluation into an auditable process, we provide a visual contrast between a standard baseline judge and our multi-step framework in Table \ref{tab:auditability_comparison}. As illustrated in Table \ref{tab:auditability_comparison}, while a baseline judge only provides an unexplainable score, the GRADE approach forces the model to externalize its criteria and justify its judgment with specific evidence from the text. This multi-step trace allows researchers to diagnose exactly why a model received a particular score, meeting the requirements for academic auditability.

\begin{table*}[htbp]
\centering
\caption{Comparison of evaluation transparency between a Baseline Judge and our GRADE Approach.}
\label{tab:auditability_comparison}
\small
\begin{tabularx}{\textwidth}{l|X|X}
\toprule
\textbf{Evaluation Stage} & \textbf{Baseline Judge (\cite{wang2024autosurvey})} & \textbf{Ours: GRADE Approach (Auditable Trace)} \\ \midrule
\textbf{1. Rubric} & Uses pre-defined generic 1--5 definitions (e.g., ``The survey has good logical consistency'') without domain context. & \textbf{[Self-Generated Rubric]} Automatically creates 7 domain-specific criteria such as: \textit{Coverage of Key ViT Axes}, \textit{Taxonomy Granularity}, etc. \\ \midrule
\textbf{2. Rubric Explanation} & \textbf{[Missing]} Provides no explanation for what it specifically measures. & \textbf{[Detailed Criterion Definition]} Explicitly defines each item, e.g., for ViT Axes: ``Measures whether the framework covers architectural variants, attention modifications, and training regimes.'' \\ \midrule
\textbf{3. Rubric Weighting} & \textbf{[Missing]}  & \textbf{[Explicit Weights]} Assigns precise importance (e.g., 0.20 for Clarity, 0.14 for Granularity). \\ \midrule
\textbf{4. Comparative Analysis} & \textbf{[Missing]} Output is often just a scalar score (e.g., \textit{Score: 3}) with no comparative grounding between articles. & \textbf{[Evidence-based Analysis]} Provides cross-article comparisons, e.g., ``Article 2 explicitly states its framework early... making it easier to apply than Article 1's implicit organization.'' \\ \midrule
\textbf{Conclusion} & \textbf{Black Box}: Human auditors cannot verify if the final score is grounded in the source text or the rubric. & \textbf{Auditable}: Human experts can inspect the rubric definitions, weights, and evidence-based analysis to validate the results. \\ \bottomrule
\end{tabularx}
\end{table*}



\section{Mapping Metrics to Bloom's Taxonomy}
\label{sec:appendix_mapping}

Table~\ref{tab:bloom_mapping} provides a systematic mapping of Bloom-Eval's 16 metrics to the six tiers of Bloom's Taxonomy. It details the theoretical alignment between standard cognitive definitions and our specific diagnostic purposes within the context of ASG.


% 定义一个垂直居中的 X 列
\newcolumntype{Y}{>{\RaggedRight\arraybackslash}X}

\begin{table*}[ht]
\centering
\caption{Systematic mapping of Bloom-Eval's 16 metrics to Bloom's Taxonomy. The table illustrates how each metric aligns with the theoretical definition of a cognitive tier to diagnose specific ASG capabilities.}
\label{tab:bloom_mapping}
\small
\renewcommand{\arraystretch}{1.6} 
\setlength{\tabcolsep}{8pt} % 增加列间距的呼吸感

\begin{tabularx}{\textwidth}{l >{\hsize=0.6\hsize}Y >{\hsize=0.7\hsize}Y >{\hsize=1.7\hsize}Y}
\toprule
\rowcolor{gray!10}
\textbf{Cognitive Level} & \textbf{Bloom's Definition} & \textbf{Bloom-Eval Metrics} & \textbf{Theoretical Alignment} \\ \midrule

L1: Memory & 
The ability to recall information, encompassing facts, concepts, definitions, and terminology. & 
$\bullet$ Entity Fidelity (EFid) \par
$\bullet$ High-Impact Reference Coverage (HIRC) \par
$\bullet$ Factual Consistency (FCons) & 
\textbf{Foundational Reproduction:} Diagnoses the system's proficiency in accurately retrieving and reproducing the field's foundational elements, specifically domain entities (EFid), pivotal literature (HIRC), and factual descriptions (FCons).  \\ \midrule

L2: Comprehension & 
The ability to explain and summarize information. & 
$\bullet$ Outline Topic Coverage (OTC) \par
$\bullet$ Citation Faithfulness (CF) \par
$\bullet$ Thematic Focus Similarity (TFSim) \par
$\bullet$ Thematic Balance (TBal) & 
\textbf{Interpretative Synthesis:} Diagnoses the ability to summarize content into necessary sub-headings (OTC), accurately explain cited references without hallucination (CF), and synthesize the research landscape's distributional hotspots (TFSim, TBal). \\ \midrule

L3: Application & 
The ability to apply learned knowledge to practical problems. & 
$\bullet$ Formatting Integrity (FMI) \par
$\bullet$ Document Structure Integrity (DSI) \par
$\bullet$ Framework Application (FAP) & 
\textbf{Standard Execution:} Diagnoses the application of standard academic survey frameworks, necessitating strict adherence to citation conventions (FMI), structural skeleton completeness (DSI), and the consistent implementation of organizing paradigms (FAP).  \\ \midrule

L4: Analysis & 
The ability to break down information and identify structures or relationships. & 
$\bullet$ Semantic Tree Similarity (STS) \par
$\bullet$ Shape Consistency (SCons) \par
$\bullet$ Structural Clarity Score (SCS) & 
\textbf{Structural Deconstruction:} Diagnoses the accurate construction of writing structure, characterized by sound hierarchical logic (STS), appropriate analytical granularity (SCons), and clear separation of distinct concepts without redundancy (SCS). \\ \midrule

L5: Evaluation & 
The ability to critically judge the value of ideas. & 
$\bullet$ Critical Analysis Alignment (CAA) & 
\textbf{Critical Reasoning:} Diagnoses the quality of evaluative reasoning by measuring the alignment of explicit value judgments (e.g., identifying method flaws) between the generated survey and expert consensus. \\ \midrule

L6: Creation & 
The capacity to propose new viewpoints, models, or works, representing the highest order of cognition. & 
$\bullet$ Framework Novelty (FNov) \par
$\bullet$ Research Outlook Quality (ROQ) & 
\textbf{Innovative Synthesis:} Diagnoses the capacity for innovation, specifically mandating the reorganization of the field with fresh conceptual models (FNov) and the foresight to identify potential future research directions (ROQ). \\ 
\bottomrule
\end{tabularx}
\end{table*}


\section{Prompt Templates}
\label{sec:appendix_prompts}

This section provides the detailed prompt templates used for the various LLM-based tasks within the Bloom-Eval framework (see Figures~\ref{fig:prompt_entity_extraction} through~\ref{fig:prompt_roq_scoring}).

\begin{figure*}[h]
    \centering
    \includegraphics[width=1\textwidth]{latex/prompt templates/level1_prompt_1.png}
    \caption{Prompt template for entity extraction.}
    \label{fig:prompt_entity_extraction}
\end{figure*}

\begin{figure*}[h]
    \centering
    \includegraphics[width=1\textwidth]{latex/prompt templates/level1_prompt_2.png}
    \caption{Prompt template for domain-adaptive entity matching.}
    \label{fig:prompt_entity_matching}
\end{figure*}

\begin{figure*}[h]
    \centering
    \includegraphics[width=1\textwidth]{latex/prompt templates/level1_prompt_3.png}
    \caption{Prompt template for factual statement extraction.}
    \label{fig:prompt_factual_extraction}
\end{figure*}

\begin{figure*}[h]
    \centering
    \includegraphics[width=1\textwidth]{latex/prompt templates/level1_prompt_4.png}
    \caption{Prompt template for factual statement matching.}
    \label{fig:prompt_factual_matching}
\end{figure*}

\begin{figure*}[h]
    \centering
    \includegraphics[width=1\textwidth]{latex/prompt templates/level2_prompt_1.png}
    \caption{Prompt template for semantic alignment of outline topics.}
    \label{fig:prompt_outline_alignment}
\end{figure*}

\begin{figure*}[h]
    \centering
    \includegraphics[width=1\textwidth]{latex/prompt templates/level5_prompt_1.png}
    \caption{Prompt template for matching critical statements.}
    \label{fig:prompt_critical_matching}
\end{figure*}

\begin{figure*}[h]
    \centering
    \includegraphics[width=1\textwidth]{latex/prompt templates/level6_prompt1.png}
    \caption{GRADE prompt for Research Outlook Quality (ROQ): Step 1(a) - rubric generation.}
    \label{fig:prompt_roq_criteria_1}
\end{figure*}

\begin{figure*}[h]
    \centering
    \includegraphics[width=1\textwidth]{latex/prompt templates/level6_prompt2.png}
    \caption{GRADE prompt for Research Outlook Quality (ROQ): Step 1(b) - rubric generation.}
    \label{fig:prompt_roq_criteria_2}
\end{figure*}

\begin{figure*}[h]
    \centering
    \includegraphics[width=1\textwidth]{latex/prompt templates/level6_prompt3.png}
    \caption{GRADE prompt for Research Outlook Quality (ROQ): Step 2 - criteria-based scoring.}
    \label{fig:prompt_roq_scoring}
\end{figure*}

\end{document}
